%%%%%%%%%%%%%%%%%%%%%%%%%%%%%%%%%%%%%%%%%%%%%%%%%%%%%%%%%%%%
%%  Real Analysis II — Course Notes (English)
%%  Part 1: Chapters 7–8
%%  From Integration to Multivariable Calculus
%%%%%%%%%%%%%%%%%%%%%%%%%%%%%%%%%%%%%%%%%%%%%%%%%%%%%%%%%%%%

\documentclass[12pt,a4paper]{book}

%% ── Encoding and fonts ──────────────────────────────────
\usepackage[utf8]{inputenc}
\usepackage[T1]{fontenc}
\usepackage{lmodern}

%% ── Page geometry ───────────────────────────────────────
\usepackage[margin=2.5cm,top=3cm,bottom=3cm]{geometry}

%% ── Mathematics ─────────────────────────────────────────
\usepackage{amsmath,amssymb,amsthm,mathtools}

%% ── Graphics and diagrams ───────────────────────────────
\usepackage{tikz}
\usetikzlibrary{calc,patterns,decorations.markings,arrows.meta}
\usepackage{pgfplots}
\pgfplotsset{compat=1.18}
\usepackage{tikz-3dplot}

%% ── Coloured boxes ──────────────────────────────────────
\usepackage[most]{tcolorbox}
\tcbuselibrary{theorems,skins,breakable}

%% ── Colours ─────────────────────────────────────────────
\usepackage{xcolor}
\definecolor{defcolor}{RGB}{0,100,180}
\definecolor{thmcolor}{RGB}{180,30,30}
\definecolor{propcolor}{RGB}{0,130,80}
\definecolor{remcolor}{RGB}{120,80,160}
\definecolor{excolor}{RGB}{200,120,0}
\definecolor{exercolor}{RGB}{60,60,60}

%% ── Lists ───────────────────────────────────────────────
\usepackage{caption}
\usepackage{enumitem}

%% ── Headers and footers ────────────────────────────────
\usepackage{fancyhdr}
\pagestyle{fancy}
\fancyhf{}
\fancyhead[LE]{\slshape\leftmark}
\fancyhead[RO]{\slshape\rightmark}
\fancyfoot[C]{\thepage}
\renewcommand{\headrulewidth}{0.4pt}

%% ── Cross-references and index ──────────────────────────
\usepackage{hyperref}
\hypersetup{
  colorlinks=true,
  linkcolor=thmcolor,
  urlcolor=defcolor,
  citecolor=propcolor
}
\usepackage{makeidx}
\makeindex

%% ════════════════════════════════════════════════════════
%%  Theorem-like environments (tcolorbox)
%% ════════════════════════════════════════════════════════

% Theorem environments (amsthm + tcolorbox wrapping)
\theoremstyle{definition}
\newtheorem{definition}{Definition}[chapter]
\newtheorem{theorem}{Theorem}[chapter]
\newtheorem{proposition}{Proposition}[chapter]
\newtheorem{lemma}{Lemma}[chapter]
\newtheorem{corollary}{Corollary}[chapter]
\newtheorem{remark}{Remark}[chapter]
\newtheorem{example}{Example}[chapter]
\newtheorem{exercise}{Exercise}[chapter]

\tcolorboxenvironment{definition}{enhanced,breakable,colback=defcolor!5,colframe=defcolor!80!black,before skip=10pt,after skip=10pt}
\tcolorboxenvironment{theorem}{enhanced,breakable,colback=thmcolor!5,colframe=thmcolor!80!black,before skip=10pt,after skip=10pt}
\tcolorboxenvironment{proposition}{enhanced,breakable,colback=propcolor!5,colframe=propcolor!80!black,before skip=10pt,after skip=10pt}
\tcolorboxenvironment{lemma}{enhanced,breakable,colback=propcolor!5,colframe=propcolor!60!black,before skip=10pt,after skip=10pt}
\tcolorboxenvironment{corollary}{enhanced,breakable,colback=thmcolor!4,colframe=thmcolor!60!black,before skip=10pt,after skip=10pt}

%% ── Plain amsthm proof environment (keep default) ──────
%% We use the standard \begin{proof}...\end{proof}

%% ════════════════════════════════════════════════════════
%%  Custom commands
%% ════════════════════════════════════════════════════════

\newcommand{\N}{\mathbb{N}}
\newcommand{\Z}{\mathbb{Z}}
\newcommand{\Q}{\mathbb{Q}}
\newcommand{\R}{\mathbb{R}}
\newcommand{\C}{\mathbb{C}}
\newcommand{\eps}{\varepsilon}
\newcommand{\abs}[1]{\left\lvert #1 \right\rvert}
\newcommand{\norm}[1]{\left\lVert #1 \right\rVert}
\newcommand{\dd}{\,\mathrm{d}}

%% ════════════════════════════════════════════════════════
%%  Title page
%% ════════════════════════════════════════════════════════

%% ════════════════════════════════════════════════════════
\begin{document}
\frontmatter

\begin{titlepage}
\begin{tikzpicture}[remember picture, overlay]
  \fill[left color=blue!15!white, right color=blue!5!white]
    (current page.south west) rectangle (current page.north east);
  \fill[color=orange!70!yellow]
    ([yshift=-2cm]current page.north west) rectangle ([yshift=-2.3cm]current page.north east);
  \fill[color=blue!80!black, rounded corners=8pt]
    ([xshift=3cm, yshift=-4cm]current page.north west) rectangle
    ([xshift=-3cm, yshift=-10cm]current page.north east);
  \node[anchor=center, text=white, font=\Huge\bfseries]
    at ([yshift=-6cm]current page.north) {Real Analysis II};
  \node[anchor=center, text=orange!80!yellow, font=\Large]
    at ([yshift=-7.5cm]current page.north) {Lecture Notes};
  \node[anchor=center, text=white!80, font=\large]
    at ([yshift=-8.8cm]current page.north) {Licence L3 --- 2025--2026};
  \node[anchor=center, text=blue!80!black, font=\Large\itshape]
    at ([yshift=-12cm]current page.north) {Ya\'e Ulrich Gaba};
  \draw[orange!70!yellow, line width=2pt]
    ([xshift=5cm, yshift=-13.5cm]current page.north west) --
    ([xshift=-5cm, yshift=-13.5cm]current page.north east);
  \node[anchor=center, text width=12cm, align=center, text=blue!60!black, font=\itshape]
    at ([yshift=-16cm]current page.north) {``One can always divide things in two: those that are continuous and those that are not.''\\[4pt]--- Henri Poincar\'e};
  \node[anchor=south, text=gray, font=\small]
    at ([yshift=1.5cm]current page.south) {\today};
\end{tikzpicture}
\end{titlepage}
\tableofcontents

%% ════════════════════════════════════════════════════════
%%  Preface
%% ════════════════════════════════════════════════════════

\chapter*{Preface}
\addcontentsline{toc}{chapter}{Preface}
\markboth{Preface}{Preface}

This volume constitutes the second semester of the first-year course in
Real Analysis.  It is a direct continuation of \emph{Real Analysis~I},
where we established the foundations of the real number system, studied
sequences and series of real numbers, explored the theory of limits and
continuity for functions of a single real variable, and developed the
differential calculus.

In the present volume we pursue three broad themes.

\begin{enumerate}[label=(\roman*)]
\item \textbf{Integration.}  We give a rigorous construction of the
  Riemann integral via Darboux sums, establish the Fundamental Theorem
  of Calculus in both of its classical forms, and develop the standard
  techniques of integration---substitution and integration by parts.  We
  also treat improper integrals and their convergence.

\item \textbf{Sequences and series of functions.}  The central concept
  here is \emph{uniform convergence}.  We prove the classical
  interchange theorems (limit with continuity, limit with integration,
  limit with differentiation), introduce the Weierstrass $M$-test, and
  apply the theory to power series.  The Weierstrass approximation
  theorem provides a beautiful conclusion to this circle of ideas.

\item \textbf{Multivariable calculus} (to appear in subsequent
  chapters).  We extend the differential and integral calculus to
  functions of several real variables, studying partial derivatives, the
  total differential, the implicit and inverse function theorems, and
  multiple integrals.
\end{enumerate}

Throughout, we maintain the same standard of rigour as in the first
semester: every major result receives a complete proof.  Numerous
examples, counterexamples, and exercises complement the theoretical
development.

\bigskip
\hfill\textit{The Authors}

%% ════════════════════════════════════════════════════════
\mainmatter
%% ════════════════════════════════════════════════════════


%%%%%%%%%%%%%%%%%%%%%%%%%%%%%%%%%%%%%%%%%%%%%%%%%%%%%%%%%%%%
%%%%%%%%%%%%%%%%%%%%%%%%%%%%%%%%%%%%%%%%%%%%%%%%%%%%%%%%%%%%
%%
%%  CHAPTER 7 : RIEMANN INTEGRATION
%%
%%%%%%%%%%%%%%%%%%%%%%%%%%%%%%%%%%%%%%%%%%%%%%%%%%%%%%%%%%%%
%%%%%%%%%%%%%%%%%%%%%%%%%%%%%%%%%%%%%%%%%%%%%%%%%%%%%%%%%%%%

\chapter{Riemann Integration}\label{ch:riemann}
\index{Riemann integral|(}

\section*{Motivation}
\addcontentsline{toc}{section}{Motivation}

The problem of computing the area of a region bounded by a curve is one
of the oldest in mathematics, going back at least to Archimedes' method
of exhaustion.  Given a non-negative continuous function
$f\colon[a,b]\to\R$, we want to assign a precise numerical value to the
``area under the curve $y=f(x)$ between $x=a$ and $x=b$.''

The basic idea is to approximate the region by rectangles.  If we divide
the interval $[a,b]$ into small subintervals and on each one erect a
rectangle whose height is the infimum (respectively the supremum) of $f$,
we obtain a \emph{lower sum} (respectively an \emph{upper sum}).  The
integral, when it exists, is the unique real number that lies between
every lower sum and every upper sum.

\begin{center}
\begin{tikzpicture}[scale=1.0]
  \begin{axis}[
    axis lines=middle,
    xlabel={$x$}, ylabel={$y$},
    xmin=-0.3, xmax=4.5,
    ymin=-0.3, ymax=3.5,
    xtick={0.5,1,1.5,2,2.5,3,3.5,4},
    xticklabels={$x_0$,$x_1$,$x_2$,$x_3$,$x_4$,$x_5$,$x_6$,$x_7$},
    ytick=\empty,
    width=12cm, height=7cm,
    every axis x label/.style={at={(ticklabel* cs:1)},anchor=west},
    every axis y label/.style={at={(ticklabel* cs:1)},anchor=south},
  ]
    % Lower sum rectangles
    \addplot[fill=blue!15, draw=blue!60, ybar interval, opacity=0.7]
      coordinates {
        (0.5,1.1) (1,1.4) (1.5,1.9) (2,2.3)
        (2.5,2.55) (3,2.6) (3.5,2.3) (4,0)
      };
    % The curve
    \addplot[domain=0.5:4, samples=100, thick, red!80!black]
      {0.3*sin(deg(1.5*x-0.5))+0.15*x+1.2};
  \end{axis}
\end{tikzpicture}
\captionof{figure}{Lower Darboux sum: rectangles with heights $m_i=\inf_{[x_{i-1},x_i]}f$.}
\end{center}

This chapter makes these ideas precise by developing the Riemann integral
in its Darboux formulation.


%% ──────────────────────────────────────────────────────
\section{Partitions and Darboux Sums}\label{sec:partitions}
\index{partition}
\index{Darboux sum}

\begin{definition}[Partition]\label{def:partition}
  Let $[a,b]$ be a closed bounded interval.  A \emph{partition}\index{partition}
  of $[a,b]$ is a finite set
  \[
    P = \{x_0, x_1, \ldots, x_n\}
    \quad\text{with}\quad
    a = x_0 < x_1 < \cdots < x_n = b.
  \]
  The intervals $[x_{i-1},x_i]$, $i=1,\ldots,n$, are called the
  \emph{subintervals} of $P$, and
  $\Delta x_i = x_i - x_{i-1}$ is the \emph{width} of the $i$-th
  subinterval.  The \emph{mesh} (or \emph{norm}) of~$P$ is
  $\norm{P} = \max_{1\le i\le n}\Delta x_i$.
\end{definition}

\begin{definition}[Darboux sums]\label{def:darboux-sums}
  Let $f\colon[a,b]\to\R$ be a bounded function and let $P=\{x_0,\ldots,x_n\}$
  be a partition of $[a,b]$.  For each $i=1,\ldots,n$ set
  \[
    m_i = \inf_{x\in[x_{i-1},x_i]} f(x),
    \qquad
    M_i = \sup_{x\in[x_{i-1},x_i]} f(x).
  \]
  The \emph{lower Darboux sum}\index{Darboux sum!lower} and
  \emph{upper Darboux sum}\index{Darboux sum!upper} of $f$ with respect
  to $P$ are
  \[
    L(f,P) = \sum_{i=1}^{n} m_i\,\Delta x_i,
    \qquad
    U(f,P) = \sum_{i=1}^{n} M_i\,\Delta x_i.
  \]
\end{definition}

\begin{center}
\begin{tikzpicture}[scale=1.0]
  \begin{axis}[
    axis lines=middle,
    xlabel={$x$}, ylabel={$y$},
    xmin=-0.3, xmax=5,
    ymin=-0.3, ymax=4,
    xtick={0.5,1.5,2.5,3.5,4.5},
    xticklabels={$x_0$,$x_1$,$x_2$,$x_3$,$x_4$},
    ytick=\empty,
    width=11cm, height=6.5cm,
    every axis x label/.style={at={(ticklabel* cs:1)},anchor=west},
    every axis y label/.style={at={(ticklabel* cs:1)},anchor=south},
    title={Upper Darboux sum (shaded) vs.\ lower Darboux sum (hatched)},
  ]
    % Upper sum rectangles
    \addplot[fill=red!15, draw=red!60, ybar interval, opacity=0.6]
      coordinates {
        (0.5,1.8) (1.5,2.6) (2.5,3.2) (3.5,3.0) (4.5,0)
      };
    % Lower sum rectangles
    \addplot[fill=blue!15, draw=blue!60, ybar interval, opacity=0.5]
      coordinates {
        (0.5,1.2) (1.5,1.8) (2.5,2.6) (3.5,2.4) (4.5,0)
      };
    % Curve
    \addplot[domain=0.5:4.5, samples=100, thick, black]
      {0.5*sin(deg(1.2*x))+0.3*x+1.0};
  \end{axis}
\end{tikzpicture}
\end{center}

\begin{definition}[Refinement]\label{def:refinement}
  A partition $P'$ is called a \emph{refinement}\index{refinement} of a
  partition $P$ if $P\subset P'$.  Given two partitions $P_1$ and $P_2$,
  their \emph{common refinement} is $P_1\cup P_2$.
\end{definition}


%% ──────────────────────────────────────────────────────
\section{Refinement Inequalities}\label{sec:refinement-ineq}

\begin{theorem}[Refinement inequalities]\label{thm:refinement-ineq}
  Let $f\colon[a,b]\to\R$ be bounded, and let $P$ and $P'$ be partitions
  of $[a,b]$ with $P\subset P'$ (i.e.\ $P'$ is a refinement of $P$).
  Then
  \[
    L(f,P) \le L(f,P') \le U(f,P') \le U(f,P).
  \]
\end{theorem}

\begin{proof}
  It suffices to prove the result when $P'$ is obtained from $P$ by
  adding a single point; the general case follows by induction.

  Let $P=\{x_0,\ldots,x_n\}$ and let $P'=P\cup\{t\}$ where
  $x_{k-1}<t<x_k$ for some index $k$.  All subintervals except
  $[x_{k-1},x_k]$ are unchanged, so we need only compare the
  contributions of $[x_{k-1},x_k]$ and its two halves
  $[x_{k-1},t]$ and $[t,x_k]$.

  Set
  \[
    m_k  = \inf_{[x_{k-1},x_k]} f,\quad
    m_k' = \inf_{[x_{k-1},t]} f,\quad
    m_k''= \inf_{[t,x_k]} f.
  \]
  Since $[x_{k-1},t]\subset[x_{k-1},x_k]$ and
  $[t,x_k]\subset[x_{k-1},x_k]$, we have $m_k\le m_k'$ and
  $m_k\le m_k''$.  Therefore
  \begin{align*}
    m_k'(t-x_{k-1}) + m_k''(x_k-t)
    &\ge m_k(t-x_{k-1}) + m_k(x_k-t) \\
    &= m_k(x_k-x_{k-1}) = m_k\,\Delta x_k.
  \end{align*}
  Hence the contribution to the lower sum does not decrease:
  $L(f,P)\le L(f,P')$.

  An entirely analogous argument, using $\sup$ and the inequality
  $M_k\ge M_k',\, M_k\ge M_k''$, shows that $U(f,P')\le U(f,P)$.

  Finally, $L(f,P')\le U(f,P')$ is immediate from $m_i\le M_i$ for
  every subinterval.
\end{proof}

\begin{corollary}[Any lower sum $\le$ any upper sum]\label{cor:lower-le-upper}
  For any two partitions $P_1,P_2$ of $[a,b]$,
  \[
    L(f,P_1) \le U(f,P_2).
  \]
\end{corollary}

\begin{proof}
  Let $P=P_1\cup P_2$.  Then $P$ refines both $P_1$ and $P_2$, so
  by Theorem~\ref{thm:refinement-ineq},
  \[
    L(f,P_1) \le L(f,P) \le U(f,P) \le U(f,P_2). \qedhere
  \]
\end{proof}


%% ──────────────────────────────────────────────────────
\section{Upper and Lower Integrals; Riemann Integrability}\label{sec:upper-lower-int}
\index{upper integral}\index{lower integral}
\index{Riemann integrability}

\begin{definition}[Upper and lower integrals]\label{def:upper-lower-int}
  Let $f\colon[a,b]\to\R$ be bounded.  The \emph{lower integral} and
  \emph{upper integral} of $f$ are
  \[
    \underline{\int_a^b} f(x)\dd x
      = \sup_P L(f,P),
    \qquad
    \overline{\int_a^b} f(x)\dd x
      = \inf_P U(f,P),
  \]
  where the supremum and infimum are taken over all partitions $P$ of
  $[a,b]$.
\end{definition}

By Corollary~\ref{cor:lower-le-upper}, every lower sum is a lower bound
for the set of upper sums, so
\[
  \underline{\int_a^b} f\dd x
  \;\le\;
  \overline{\int_a^b} f\dd x.
\]

\begin{definition}[Riemann integrability]\label{def:riemann-int}
  A bounded function $f\colon[a,b]\to\R$ is said to be
  \emph{Riemann integrable}\index{Riemann integral!integrability} on
  $[a,b]$ if
  \[
    \underline{\int_a^b} f(x)\dd x
    = \overline{\int_a^b} f(x)\dd x.
  \]
  In that case, the common value is called the \emph{Riemann integral}
  of $f$ on $[a,b]$ and is denoted
  \[
    \int_a^b f(x)\dd x.
  \]
  We write $\mathcal{R}([a,b])$ for the set of Riemann-integrable
  functions on $[a,b]$.
\end{definition}


%% ──────────────────────────────────────────────────────
\section{Integrability Criterion}\label{sec:int-criterion}

The following criterion is the workhorse for proving integrability.

\begin{theorem}[Riemann integrability criterion]\label{thm:int-criterion}
  A bounded function $f\colon[a,b]\to\R$ is Riemann integrable if and
  only if for every $\eps>0$ there exists a partition $P$ of $[a,b]$ such
  that
  \[
    U(f,P) - L(f,P) < \eps.
  \]
\end{theorem}

\begin{proof}
  \textbf{($\Rightarrow$).}\;
  Suppose $f\in\mathcal{R}([a,b])$ and let $I=\int_a^b f\dd x$.
  Given $\eps>0$, by definition of supremum and infimum there exist
  partitions $P_1,P_2$ such that
  \[
    I - \frac{\eps}{2} < L(f,P_1)
    \quad\text{and}\quad
    U(f,P_2) < I + \frac{\eps}{2}.
  \]
  Let $P=P_1\cup P_2$.  By the refinement inequalities,
  \[
    U(f,P) - L(f,P) \le U(f,P_2) - L(f,P_1)
    < \Bigl(I+\frac{\eps}{2}\Bigr) - \Bigl(I-\frac{\eps}{2}\Bigr) = \eps.
  \]

  \textbf{($\Leftarrow$).}\;
  Conversely, suppose that for every $\eps>0$ there exists a partition $P$
  with $U(f,P)-L(f,P)<\eps$.  Since
  \[
    0 \le \overline{\int_a^b}f\dd x - \underline{\int_a^b}f\dd x
    \le U(f,P) - L(f,P) < \eps
  \]
  for every $\eps>0$, we conclude that the upper and lower integrals are
  equal, so $f$ is integrable.
\end{proof}


%% ──────────────────────────────────────────────────────
\section{Continuous Functions Are Integrable}\label{sec:cont-integrable}

\begin{theorem}[Integrability of continuous functions]\label{thm:cont-integrable}
  \index{Riemann integral!continuous functions}
  Every continuous function $f\colon[a,b]\to\R$ is Riemann integrable.
\end{theorem}

\begin{proof}
  Since $[a,b]$ is compact, $f$ is \emph{uniformly} continuous on
  $[a,b]$: given $\eps>0$, there exists $\delta>0$ such that
  \[
    \abs{x-y}<\delta \implies \abs{f(x)-f(y)}<\frac{\eps}{b-a}.
  \]
  Choose a partition $P=\{x_0,\ldots,x_n\}$ of $[a,b]$ with
  $\norm{P}<\delta$.  On each subinterval $[x_{i-1},x_i]$ the function
  $f$ attains its supremum $M_i$ and infimum $m_i$ (by the extreme
  value theorem), say at points $u_i$ and $v_i$ respectively.  Since
  $\abs{u_i-v_i}\le\Delta x_i<\delta$, we have
  \[
    M_i - m_i = f(u_i)-f(v_i) < \frac{\eps}{b-a}.
  \]
  Therefore
  \[
    U(f,P) - L(f,P)
    = \sum_{i=1}^n (M_i-m_i)\,\Delta x_i
    < \frac{\eps}{b-a}\sum_{i=1}^n \Delta x_i
    = \frac{\eps}{b-a}\cdot(b-a) = \eps.
  \]
  By Theorem~\ref{thm:int-criterion}, $f$ is Riemann integrable.
\end{proof}


%% ──────────────────────────────────────────────────────
\section{Monotone Functions Are Integrable}\label{sec:mono-integrable}

\begin{theorem}[Integrability of monotone functions]\label{thm:mono-integrable}
  \index{Riemann integral!monotone functions}
  Every monotone function $f\colon[a,b]\to\R$ is Riemann integrable.
\end{theorem}

\begin{proof}
  Assume $f$ is non-decreasing (the non-increasing case is analogous).
  Then $f$ is bounded, with $f(a)\le f(x)\le f(b)$ for all $x\in[a,b]$.
  On each subinterval $[x_{i-1},x_i]$, we have
  $m_i=f(x_{i-1})$ and $M_i=f(x_i)$.

  Given $\eps>0$, choose the uniform partition $P$ with $n$ subintervals
  of equal width $\Delta x_i=(b-a)/n$.  Then
  \begin{align*}
    U(f,P)-L(f,P)
    &= \sum_{i=1}^n \bigl(f(x_i)-f(x_{i-1})\bigr)\,\frac{b-a}{n} \\
    &= \frac{b-a}{n}\sum_{i=1}^n \bigl(f(x_i)-f(x_{i-1})\bigr) \\
    &= \frac{b-a}{n}\bigl(f(b)-f(a)\bigr).
  \end{align*}
  The last expression is a telescoping sum.  Choosing $n$ large enough so
  that
  \[
    \frac{b-a}{n}\bigl(f(b)-f(a)\bigr) < \eps,
  \]
  i.e.\ $n > (b-a)(f(b)-f(a))/\eps$, we conclude by the integrability
  criterion.
\end{proof}


%% ──────────────────────────────────────────────────────
\section{Piecewise Continuous Functions}\label{sec:piecewise}
\index{piecewise continuous function}

\begin{definition}[Piecewise continuous function]\label{def:piecewise}
  A function $f\colon[a,b]\to\R$ is \emph{piecewise continuous} if there
  exists a partition $a=c_0<c_1<\cdots<c_m=b$ such that $f$ is
  continuous on each open subinterval $(c_{j-1},c_j)$ and the one-sided
  limits $\lim_{x\to c_j^+}f(x)$ and $\lim_{x\to c_j^-}f(x)$ exist
  (and are finite) at each partition point.
\end{definition}

\begin{proposition}[Integrability of piecewise continuous functions]\label{prop:piecewise-int}
  Every piecewise continuous function on $[a,b]$ is Riemann integrable.
\end{proposition}

\begin{proof}
  On each subinterval $[c_{j-1},c_j]$, the function $f$ is bounded
  (since the one-sided limits exist) and has at most finitely many
  discontinuities (at the endpoints).  One shows, using the integrability
  criterion, that a bounded function with finitely many discontinuities
  is integrable.  Indeed, near each discontinuity one can choose
  subintervals of total width less than $\eps/(4M)$ (where
  $M=\sup\abs{f}$), and on the remaining subintervals the uniform
  continuity argument from Theorem~\ref{thm:cont-integrable} applies.
  By the Chasles relation (Proposition~\ref{prop:chasles}), the integral
  over $[a,b]$ is the sum of the integrals over the subintervals.
\end{proof}


%% ──────────────────────────────────────────────────────
\section{Properties of the Riemann Integral}\label{sec:int-properties}
\index{Riemann integral!properties}

\begin{proposition}[Linearity]\label{prop:linearity}
  If $f,g\in\mathcal{R}([a,b])$ and $\alpha,\beta\in\R$, then
  $\alpha f+\beta g\in\mathcal{R}([a,b])$ and
  \[
    \int_a^b\bigl(\alpha f(x)+\beta g(x)\bigr)\dd x
    = \alpha\int_a^b f(x)\dd x + \beta\int_a^b g(x)\dd x.
  \]
\end{proposition}

\begin{proof}
  We prove the result for $\alpha f$ and $f+g$ separately.

  \emph{Scalar multiple.}\;
  Suppose $\alpha\ge0$ (the case $\alpha<0$ is similar, exchanging sup
  and inf).  For any partition $P$ and any subinterval $[x_{i-1},x_i]$,
  \[
    \inf_{[x_{i-1},x_i]}(\alpha f) = \alpha\inf_{[x_{i-1},x_i]}f,
    \qquad
    \sup_{[x_{i-1},x_i]}(\alpha f) = \alpha\sup_{[x_{i-1},x_i]}f.
  \]
  Hence $L(\alpha f,P)=\alpha L(f,P)$ and $U(\alpha f,P)=\alpha U(f,P)$.
  Taking supremum and infimum over all $P$ gives the result.

  \emph{Sum.}\;
  For any subinterval,
  \[
    \inf(f+g) \ge \inf f + \inf g
    \quad\text{and}\quad
    \sup(f+g) \le \sup f + \sup g.
  \]
  Therefore $L(f+g,P)\ge L(f,P)+L(g,P)$ and
  $U(f+g,P)\le U(f,P)+U(g,P)$.  It follows that
  \begin{align*}
    U(f+g,P)-L(f+g,P)
    &\le \bigl(U(f,P)-L(f,P)\bigr)+\bigl(U(g,P)-L(g,P)\bigr).
  \end{align*}
  Given $\eps>0$, choose $P_1$ with $U(f,P_1)-L(f,P_1)<\eps/2$ and
  $P_2$ with $U(g,P_2)-L(g,P_2)<\eps/2$; set $P=P_1\cup P_2$.  Then
  \[
    U(f+g,P)-L(f+g,P) < \eps,
  \]
  so $f+g$ is integrable.  The value of the integral follows from
  \[
    L(f,P)+L(g,P) \le L(f+g,P) \le \int_a^b(f+g) \le U(f+g,P)
    \le U(f,P)+U(g,P),
  \]
  letting the partition refine and noting that the outer bounds converge
  to $\int f + \int g$.
\end{proof}

\begin{proposition}[Positivity and monotonicity]\label{prop:positivity}
  \index{Riemann integral!positivity}
  Let $f,g\in\mathcal{R}([a,b])$.
  \begin{enumerate}[label=(\alph*)]
    \item \textbf{Positivity.}\; If $f(x)\ge0$ for all $x\in[a,b]$, then
      $\displaystyle\int_a^b f(x)\dd x\ge0$.
    \item \textbf{Monotonicity.}\; If $f(x)\le g(x)$ for all $x\in[a,b]$,
      then $\displaystyle\int_a^b f(x)\dd x\le\int_a^b g(x)\dd x$.
  \end{enumerate}
\end{proposition}

\begin{proof}
  (a) If $f\ge0$ then $m_i\ge0$ for every subinterval, so
  $L(f,P)\ge0$ for every partition $P$.  Hence
  $\int_a^b f\dd x=\sup_P L(f,P)\ge0$.

  (b) Apply~(a) to $g-f\ge0$ and use linearity.
\end{proof}

\begin{proposition}[Chasles relation]\label{prop:chasles}
  \index{Chasles relation}
  If $f\in\mathcal{R}([a,b])$ and $a<c<b$, then
  $f\in\mathcal{R}([a,c])\cap\mathcal{R}([c,b])$ and
  \[
    \int_a^b f(x)\dd x = \int_a^c f(x)\dd x + \int_c^b f(x)\dd x.
  \]
\end{proposition}

\begin{proof}
  Given $\eps>0$, choose a partition $P$ of $[a,b]$ with
  $U(f,P)-L(f,P)<\eps$.  Let $P'=P\cup\{c\}$; then
  $U(f,P')-L(f,P')\le U(f,P)-L(f,P)<\eps$ by the refinement inequality.
  Write $P'=P_1\cup P_2$ where $P_1=P'\cap[a,c]$ is a partition of
  $[a,c]$ and $P_2=P'\cap[c,b]$ is a partition of $[c,b]$.  Then
  \begin{align*}
    U(f,P_1)-L(f,P_1) + U(f,P_2)-L(f,P_2)
    &= U(f,P')-L(f,P') < \eps.
  \end{align*}
  Since each term on the left is non-negative, both are less than $\eps$.
  Hence $f$ is integrable on $[a,c]$ and $[c,b]$.  The additivity of the
  integral values follows from $L(f,P')=L(f,P_1)+L(f,P_2)$ and similarly
  for upper sums.
\end{proof}


%% ──────────────────────────────────────────────────────
\section{The Triangle Inequality for Integrals}\label{sec:int-triangle}

\begin{theorem}[Mean value inequality]\label{thm:mean-value-ineq}
  \index{integral!triangle inequality}
  If $f\in\mathcal{R}([a,b])$ then $\abs{f}\in\mathcal{R}([a,b])$ and
  \[
    \abs*{\int_a^b f(x)\dd x} \le \int_a^b \abs{f(x)}\dd x.
  \]
\end{theorem}

\begin{proof}
  First we show $\abs{f}$ is integrable.  For any subinterval
  $[x_{i-1},x_i]$,
  \[
    \sup_{[x_{i-1},x_i]}\abs{f} - \inf_{[x_{i-1},x_i]}\abs{f}
    \le \sup_{[x_{i-1},x_i]} f - \inf_{[x_{i-1},x_i]} f,
  \]
  since $\bigl|\abs{f(x)}-\abs{f(y)}\bigr|\le\abs{f(x)-f(y)}$.
  Therefore $U(\abs{f},P)-L(\abs{f},P)\le U(f,P)-L(f,P)$ for every
  partition $P$, and integrability of $\abs{f}$ follows from the
  integrability criterion.

  For the inequality, note that $-\abs{f}\le f\le\abs{f}$.  By
  monotonicity of the integral,
  \[
    -\int_a^b\abs{f}\dd x \le \int_a^b f\dd x \le \int_a^b\abs{f}\dd x,
  \]
  which is precisely $\abs{\int_a^b f\dd x}\le\int_a^b\abs{f}\dd x$.
\end{proof}


%% ──────────────────────────────────────────────────────
\section{The Fundamental Theorem of Calculus}\label{sec:FTC}
\index{Fundamental Theorem of Calculus|(}

The Fundamental Theorem of Calculus is one of the most important results
in all of analysis.  It establishes the deep connection between the
differential and integral calculus.

\subsection{First Fundamental Theorem}\label{subsec:FTC1}

\begin{theorem}[First Fundamental Theorem of Calculus]\label{thm:FTC1}
  Let $f\colon[a,b]\to\R$ be Riemann integrable, and define the
  \emph{area function}\index{area function}
  \[
    F(x) = \int_a^x f(t)\dd t, \qquad x\in[a,b].
  \]
  Then:
  \begin{enumerate}[label=(\alph*)]
    \item $F$ is continuous on $[a,b]$.
    \item If $f$ is continuous at a point $c\in(a,b)$, then $F$ is
      differentiable at $c$ and $F'(c)=f(c)$.
  \end{enumerate}
  In particular, if $f$ is continuous on $[a,b]$, then $F$ is a
  $C^1$ antiderivative of $f$ on $[a,b]$.
\end{theorem}

\begin{proof}
  \textbf{(a) Continuity of $F$.}\;
  Since $f$ is Riemann integrable on $[a,b]$, it is bounded; let
  $M=\sup_{[a,b]}\abs{f}$.  For $x,y\in[a,b]$ with $x<y$,
  \[
    \abs{F(y)-F(x)} = \abs*{\int_x^y f(t)\dd t}
    \le \int_x^y \abs{f(t)}\dd t \le M(y-x).
  \]
  Given $\eps>0$, choosing $\delta=\eps/M$ (or $\delta=1$ if $M=0$)
  shows that $\abs{y-x}<\delta$ implies $\abs{F(y)-F(x)}<\eps$.

  \textbf{(b) Differentiability at a point of continuity.}\;
  Let $c\in(a,b)$ and suppose $f$ is continuous at $c$.  For $h\ne0$
  with $c+h\in[a,b]$,
  \[
    \frac{F(c+h)-F(c)}{h}
    = \frac{1}{h}\int_c^{c+h} f(t)\dd t.
  \]
  We must show this converges to $f(c)$ as $h\to0$.  Write
  \[
    \frac{1}{h}\int_c^{c+h} f(t)\dd t - f(c)
    = \frac{1}{h}\int_c^{c+h}\bigl(f(t)-f(c)\bigr)\dd t.
  \]
  Given $\eps>0$, by continuity of $f$ at $c$ there exists $\delta>0$
  such that $\abs{t-c}<\delta$ implies $\abs{f(t)-f(c)}<\eps$.  For
  $0<\abs{h}<\delta$, every $t$ between $c$ and $c+h$ satisfies
  $\abs{t-c}\le\abs{h}<\delta$, so
  \[
    \abs*{\frac{1}{h}\int_c^{c+h}\bigl(f(t)-f(c)\bigr)\dd t}
    \le \frac{1}{\abs{h}}\int_{\min(c,c+h)}^{\max(c,c+h)}
       \abs{f(t)-f(c)}\dd t
    < \frac{1}{\abs{h}}\cdot\eps\cdot\abs{h} = \eps.
  \]
  Therefore $F'(c)=f(c)$.
\end{proof}

\begin{center}
\begin{tikzpicture}
  \begin{axis}[
    axis lines=middle,
    xlabel={$x$}, ylabel={},
    xmin=-0.5, xmax=5,
    ymin=-0.5, ymax=5,
    xtick={0.5,3.5},
    xticklabels={$a$,$x$},
    ytick=\empty,
    width=11cm, height=6cm,
    every axis x label/.style={at={(ticklabel* cs:1)},anchor=west},
  ]
    % Shaded area
    \addplot[fill=blue!15, draw=none, domain=0.5:3.5, samples=100]
      {0.3*sin(deg(1.5*x))+0.15*x+1.5} \closedcycle;
    % Curve f
    \addplot[domain=0.2:4.8, samples=100, thick, red!80!black]
      {0.3*sin(deg(1.5*x))+0.15*x+1.5};
    \node at (axis cs:4.2,2.6) {$y=f(t)$};
    \node at (axis cs:2,0.8) {$F(x)=\displaystyle\int_a^x f(t)\dd t$};
    % Vertical lines
    \addplot[dashed, thin, black] coordinates {(0.5,0) (0.5,1.72)};
    \addplot[dashed, thin, black] coordinates {(3.5,0) (3.5,2.44)};
  \end{axis}
\end{tikzpicture}
\end{center}


\subsection{Second Fundamental Theorem (Newton--Leibniz formula)}\label{subsec:FTC2}

\begin{theorem}[Second Fundamental Theorem of Calculus]\label{thm:FTC2}
  \index{Newton--Leibniz formula}
  Let $f\colon[a,b]\to\R$ be Riemann integrable.  If $G\colon[a,b]\to\R$
  is continuous on $[a,b]$, differentiable on $(a,b)$, and $G'(x)=f(x)$
  for all $x\in(a,b)$, then
  \[
    \int_a^b f(x)\dd x = G(b)-G(a).
  \]
\end{theorem}

\begin{proof}
  Let $P=\{x_0,\ldots,x_n\}$ be any partition of $[a,b]$.  By the Mean
  Value Theorem, for each $i$ there exists $c_i\in(x_{i-1},x_i)$ with
  \[
    G(x_i)-G(x_{i-1}) = G'(c_i)\,\Delta x_i = f(c_i)\,\Delta x_i.
  \]
  Summing over $i$ and using the telescoping property,
  \[
    G(b)-G(a) = \sum_{i=1}^n f(c_i)\,\Delta x_i.
  \]
  Since $m_i\le f(c_i)\le M_i$, we have
  \[
    L(f,P) \le G(b)-G(a) \le U(f,P).
  \]
  This holds for every partition $P$.  Because $f$ is integrable,
  $\sup_P L(f,P)=\inf_P U(f,P)=\int_a^b f\dd x$, and the only number
  that lies between all lower and all upper sums is the integral.
  Therefore $G(b)-G(a)=\int_a^b f\dd x$.
\end{proof}

\begin{remark}[Notation for antiderivative evaluation]\label{rem:antideriv-notation}
  We write $\bigl[G(x)\bigr]_a^b = G(b)-G(a)$.  The Second Fundamental
  Theorem thus reads
  $\displaystyle\int_a^b G'(x)\dd x = \bigl[G(x)\bigr]_a^b$.
\end{remark}

\index{Fundamental Theorem of Calculus|)}


%% ──────────────────────────────────────────────────────
\section{Integration Techniques}\label{sec:integration-techniques}

\subsection{Integration by Parts}\label{subsec:int-by-parts}
\index{integration by parts}

\begin{theorem}[Integration by parts]\label{thm:int-by-parts}
  Let $u,v\colon[a,b]\to\R$ be continuously differentiable.  Then
  \[
    \int_a^b u(x)\,v'(x)\dd x
    = \bigl[u(x)\,v(x)\bigr]_a^b - \int_a^b u'(x)\,v(x)\dd x.
  \]
\end{theorem}

\begin{proof}
  By the product rule, $(uv)'=u'v+uv'$, so $uv'=(uv)'-u'v$.
  Since $u$ and $v$ are $C^1$, all three functions $uv'$, $(uv)'$, and
  $u'v$ are continuous, hence integrable.  Integrating from $a$ to $b$
  and applying the Second Fundamental Theorem to $(uv)'$,
  \[
    \int_a^b u\,v'\dd x
    = \int_a^b (uv)'\dd x - \int_a^b u'v\dd x
    = \bigl[uv\bigr]_a^b - \int_a^b u'v\dd x. \qedhere
  \]
\end{proof}

\begin{example}[Integration by parts]\label{ex:ibp-example1}
  Compute $\displaystyle I=\int_0^1 x\,e^x\dd x$.

  Set $u(x)=x$, $v'(x)=e^x$, so $u'(x)=1$, $v(x)=e^x$.  Then
  \[
    I = \bigl[x\,e^x\bigr]_0^1 - \int_0^1 e^x\dd x
    = e - \bigl[e^x\bigr]_0^1
    = e-(e-1) = 1.
  \]
\end{example}

\begin{example}[Reduction formula]\label{ex:ibp-reduction}
  For $n\ge1$, let $I_n=\int_0^{\pi/2}\sin^n x\dd x$.  Integration by
  parts with $u=\sin^{n-1}x$, $v'=\sin x$ gives
  \[
    I_n = \frac{n-1}{n}\,I_{n-2} \qquad(n\ge2),
  \]
  with $I_0=\pi/2$ and $I_1=1$.
\end{example}


\subsection{Integration by Substitution}\label{subsec:substitution}
\index{integration by substitution}

\begin{theorem}[Change of variable]\label{thm:substitution}
  Let $\varphi\colon[\alpha,\beta]\to\R$ be a $C^1$ function with
  $\varphi([\alpha,\beta])\subset[a,b]$, and let $f\colon[a,b]\to\R$ be
  continuous.  Then
  \[
    \int_{\varphi(\alpha)}^{\varphi(\beta)} f(x)\dd x
    = \int_\alpha^\beta f\bigl(\varphi(t)\bigr)\,\varphi'(t)\dd t.
  \]
\end{theorem}

\begin{proof}
  Let $F$ be an antiderivative of $f$ (which exists by the First
  Fundamental Theorem, since $f$ is continuous).  Define
  $H(t)=F(\varphi(t))$.  By the chain rule,
  \[
    H'(t) = F'(\varphi(t))\,\varphi'(t) = f(\varphi(t))\,\varphi'(t).
  \]
  Since $H'$ is continuous, the Second Fundamental Theorem gives
  \[
    \int_\alpha^\beta f(\varphi(t))\,\varphi'(t)\dd t
    = H(\beta)-H(\alpha)
    = F(\varphi(\beta))-F(\varphi(\alpha))
    = \int_{\varphi(\alpha)}^{\varphi(\beta)} f(x)\dd x. \qedhere
  \]
\end{proof}

\begin{example}[Substitution example]\label{ex:subst-example1}
  Compute $\displaystyle\int_0^1 \frac{2x}{1+x^2}\dd x$.

  Set $\varphi(t)=t$, $x=t$, or more directly, let $u=1+x^2$, so
  $\dd u=2x\dd x$.  When $x=0$, $u=1$; when $x=1$, $u=2$.  Thus
  \[
    \int_0^1 \frac{2x}{1+x^2}\dd x
    = \int_1^2 \frac{\dd u}{u}
    = \bigl[\ln u\bigr]_1^2 = \ln 2.
  \]
\end{example}

\begin{example}[Trigonometric substitution]\label{ex:subst-trig}
  Compute $\displaystyle\int_0^1 \sqrt{1-x^2}\dd x$.

  Set $x=\sin\theta$, $\dd x=\cos\theta\dd\theta$.  When $x=0$,
  $\theta=0$; when $x=1$, $\theta=\pi/2$.  Then
  \[
    \int_0^1\sqrt{1-x^2}\dd x
    = \int_0^{\pi/2}\cos^2\theta\dd\theta
    = \frac{\pi}{4}.
  \]
  This confirms the area of a quarter-disk of radius~$1$.
\end{example}


%% ──────────────────────────────────────────────────────
\section{Improper Integrals}\label{sec:improper}
\index{improper integral|(}

\begin{definition}[Improper integral on $[a,+\infty)$]\label{def:improper-inf}
  Let $f\colon[a,+\infty)\to\R$ be Riemann integrable on $[a,T]$ for
  every $T>a$.  If the limit
  \[
    \lim_{T\to+\infty}\int_a^T f(x)\dd x
  \]
  exists and is finite, we say the improper integral
  $\displaystyle\int_a^{+\infty}f(x)\dd x$ \emph{converges} and we set
  \[
    \int_a^{+\infty}f(x)\dd x = \lim_{T\to+\infty}\int_a^T f(x)\dd x.
  \]
  Otherwise, we say the integral \emph{diverges}.
\end{definition}

\begin{definition}[Improper integral at a singularity]\label{def:improper-sing}
  Let $f\colon(a,b]\to\R$ be Riemann integrable on $[a+\eta,b]$ for
  every $0<\eta<b-a$.  If $\displaystyle\lim_{\eta\to0^+}\int_{a+\eta}^b f(x)\dd x$
  exists and is finite, we set
  \[
    \int_a^b f(x)\dd x = \lim_{\eta\to0^+}\int_{a+\eta}^b f(x)\dd x.
  \]
\end{definition}

\begin{example}{The integral $\int_1^{+\infty}x^{-\alpha}\dd x$}{riemann-alpha}
  \index{Riemann integral!$p$-integral}
  For $\alpha\ne1$:
  \[
    \int_1^T x^{-\alpha}\dd x
    = \left[\frac{x^{1-\alpha}}{1-\alpha}\right]_1^T
    = \frac{T^{1-\alpha}-1}{1-\alpha}.
  \]
  \begin{itemize}
    \item If $\alpha>1$: $T^{1-\alpha}\to0$ as $T\to+\infty$, so the
      integral converges to $\dfrac{1}{\alpha-1}$.
    \item If $\alpha<1$: $T^{1-\alpha}\to+\infty$, so the integral
      diverges.
    \item If $\alpha=1$: $\displaystyle\int_1^T\frac{\dd x}{x}=\ln T\to+\infty$,
      so the integral diverges.
  \end{itemize}
  Conclusion: $\displaystyle\int_1^{+\infty}\frac{\dd x}{x^\alpha}$
  converges if and only if $\alpha>1$.
\end{example}

\begin{proposition}[Comparison test for improper integrals]\label{prop:comparison-test}
  \index{improper integral!comparison test}
  Let $0\le f(x)\le g(x)$ for all $x\ge a$.
  \begin{enumerate}[label=(\alph*)]
    \item If $\int_a^{+\infty}g\dd x$ converges, then so does
      $\int_a^{+\infty}f\dd x$.
    \item If $\int_a^{+\infty}f\dd x$ diverges, then so does
      $\int_a^{+\infty}g\dd x$.
  \end{enumerate}
\end{proposition}

\begin{proof}
  (a) The function $F(T)=\int_a^T f\dd x$ is non-decreasing and bounded
  above by $\int_a^{+\infty}g\dd x$, hence has a finite limit.

  (b) Contrapositive of~(a).
\end{proof}

\begin{definition}[Absolute convergence]\label{def:abs-conv}
  An improper integral $\int_a^{+\infty}f\dd x$ is
  \emph{absolutely convergent}\index{improper integral!absolute convergence}
  if $\int_a^{+\infty}\abs{f}\dd x$ converges.
\end{definition}

\begin{proposition}[Absolute convergence implies convergence]\label{prop:abs-implies-conv}
  If $\int_a^{+\infty}\abs{f(x)}\dd x$ converges, then
  $\int_a^{+\infty}f(x)\dd x$ converges.
\end{proposition}

\begin{proof}
  Write $f=f^+-f^-$ where $f^+=\max(f,0)$ and $f^-=\max(-f,0)$.  Then
  $0\le f^\pm\le\abs{f}$, so both $\int f^+$ and $\int f^-$ converge
  by the comparison test.  Hence $\int f=\int f^+-\int f^-$ converges.
\end{proof}

\index{improper integral|)}


%% ──────────────────────────────────────────────────────
\section{Wallis and Stirling Formulas}\label{sec:wallis-stirling}

\begin{theorem}[Wallis formula]\label{thm:wallis}
  \index{Wallis formula}
  \[
    \lim_{n\to\infty}
    \frac{1}{2n+1}
    \left(\frac{(2n)!!}{(2n-1)!!}\right)^2
    = \frac{\pi}{2},
  \]
  or equivalently,
  \[
    \frac{\pi}{2}
    = \prod_{n=1}^{\infty}\frac{4n^2}{4n^2-1}
    = \frac{2}{1}\cdot\frac{2}{3}\cdot\frac{4}{3}\cdot\frac{4}{5}
      \cdot\frac{6}{5}\cdot\frac{6}{7}\cdots
  \]
\end{theorem}

\begin{theorem}[Stirling formula]\label{thm:stirling}
  \index{Stirling formula}
  As $n\to\infty$,
  \[
    n! \sim \sqrt{2\pi n}\left(\frac{n}{e}\right)^n,
  \]
  meaning $\displaystyle\lim_{n\to\infty}
  \frac{n!}{\sqrt{2\pi n}\,(n/e)^n}=1$.
\end{theorem}

These classical formulas connect the Wallis integrals $I_n=\int_0^{\pi/2}\sin^n x\dd x$
with the constants $\pi$ and $e$.  Their proofs, which rely on the
reduction formula of Example~\ref{ex:ibp-reduction} and careful
asymptotic analysis, can be found in more advanced treatments.


%% ──────────────────────────────────────────────────────
\section{Exercises}\label{sec:exer-ch7}

\begin{exercise}[Direct computation via Darboux sums]\label{exo:exer-darboux-direct}
  Let $f(x)=x^2$ on $[0,1]$.  Using the uniform partition with $n$
  subintervals, compute $L(f,P_n)$ and $U(f,P_n)$ explicitly.  Verify
  that both converge to $1/3$ as $n\to\infty$, and conclude that
  $\int_0^1 x^2\dd x=1/3$.
\end{exercise}

\begin{exercise}[Non-integrable function]\label{exo:exer-dirichlet}
  Show that the Dirichlet function $\mathbf{1}_\Q$ (equal to~$1$ on~$\Q$
  and~$0$ on~$\R\setminus\Q$) is not Riemann integrable on any interval
  $[a,b]$ with $a<b$.
\end{exercise}

\begin{exercise}[Integrability of $f^2$]\label{exo:exer-f-squared}
  Prove that if $f\in\mathcal{R}([a,b])$, then
  $f^2\in\mathcal{R}([a,b])$.
  \emph{Hint:} use $M_i^2-m_i^2=(M_i+m_i)(M_i-m_i)\le 2M(M_i-m_i)$.
\end{exercise}

\begin{exercise}[Product of integrable functions]\label{exo:exer-product}
  Prove that if $f,g\in\mathcal{R}([a,b])$, then
  $fg\in\mathcal{R}([a,b])$.
  \emph{Hint:} write $fg=\tfrac{1}{4}\bigl((f+g)^2-(f-g)^2\bigr)$.
\end{exercise}

\begin{exercise}[Mean Value Theorem for integrals]\label{exo:exer-MVT-int}
  Let $f\colon[a,b]\to\R$ be continuous.  Prove that there exists
  $c\in[a,b]$ such that
  $\displaystyle\int_a^b f(x)\dd x = f(c)(b-a)$.
\end{exercise}

\begin{exercise}[Second Mean Value Theorem]\label{exo:exer-second-MVT}
  Let $f\colon[a,b]\to\R$ be continuous and $g\colon[a,b]\to\R$ be
  integrable with $g\ge0$.  Prove that there exists $c\in[a,b]$ such
  that $\displaystyle\int_a^b f(x)g(x)\dd x = f(c)\int_a^b g(x)\dd x$.
\end{exercise}

\begin{exercise}[Integration by parts]\label{exo:exer-ibp}
  Compute the following integrals using integration by parts:
  \begin{enumerate}[label=(\alph*)]
    \item $\displaystyle\int_0^1 x^2 e^x\dd x$
    \item $\displaystyle\int_0^{\pi} x\sin x\dd x$
    \item $\displaystyle\int_1^e (\ln x)^2\dd x$
  \end{enumerate}
\end{exercise}

\begin{exercise}[Substitution practice]\label{exo:exer-subst}
  Evaluate:
  \begin{enumerate}[label=(\alph*)]
    \item $\displaystyle\int_0^{\pi/4}\tan x\dd x$
    \item $\displaystyle\int_0^1 \frac{x}{(1+x^2)^2}\dd x$
    \item $\displaystyle\int_0^{+\infty} e^{-x}\sin x\dd x$
  \end{enumerate}
\end{exercise}

\begin{exercise}[Convergence of improper integrals]\label{exo:exer-improper-conv}
  Determine whether each integral converges or diverges:
  \begin{enumerate}[label=(\alph*)]
    \item $\displaystyle\int_1^{+\infty}\frac{\dd x}{x^2+1}$
    \item $\displaystyle\int_0^1\frac{\dd x}{\sqrt{x}}$
    \item $\displaystyle\int_1^{+\infty}\frac{\sin x}{x^2}\dd x$
    \item $\displaystyle\int_0^{+\infty}\frac{\dd x}{1+x^3}$
  \end{enumerate}
\end{exercise}

\begin{exercise}[Integral remainder for Taylor series]\label{exo:exer-taylor-remainder}
  Let $f\in C^{n+1}([a,b])$.  Using integration by parts $n$ times,
  derive the integral form of the Taylor remainder:
  \[
    f(b)=\sum_{k=0}^n \frac{f^{(k)}(a)}{k!}(b-a)^k
    + \int_a^b \frac{f^{(n+1)}(t)}{n!}(b-t)^n\dd t.
  \]
\end{exercise}

\begin{exercise}[Cauchy--Schwarz inequality for integrals]\label{exo:exer-CS}
  Let $f,g\in\mathcal{R}([a,b])$.  Prove that
  \[
    \left(\int_a^b f(x)g(x)\dd x\right)^2
    \le \int_a^b f(x)^2\dd x\;\cdot\;\int_a^b g(x)^2\dd x.
  \]
  \emph{Hint:} for $\lambda\in\R$, expand
  $\int_a^b(f+\lambda g)^2\dd x\ge0$.
\end{exercise}

\begin{exercise}[An integral identity]\label{exo:exer-integral-identity}
  Let $f\colon[0,\pi]\to\R$ be continuous.  Prove that
  \[
    \int_0^\pi x\,f(\sin x)\dd x = \frac{\pi}{2}\int_0^\pi f(\sin x)\dd x.
  \]
  \emph{Hint:} use the substitution $u=\pi-x$.
\end{exercise}

\bigskip
\begin{tcolorbox}[colback=gray!5, colframe=gray!60, title={\textbf{Chapter 7 --- Summary}}]
\begin{itemize}[leftmargin=*]
  \item A bounded function $f\colon[a,b]\to\R$ is \textbf{Riemann integrable}
    when its upper and lower Darboux integrals coincide.
  \item The \textbf{integrability criterion}: $f$ is integrable iff
    $\forall\eps>0$, $\exists P$ with $U(f,P)-L(f,P)<\eps$.
  \item \textbf{Continuous} and \textbf{monotone} functions on $[a,b]$ are
    integrable.
  \item The integral is \textbf{linear}, \textbf{monotone}, and satisfies the
    \textbf{Chasles relation} and the \textbf{triangle inequality}.
  \item The \textbf{First FTC}: if $f$ is continuous, then
    $F(x)=\int_a^x f$ is differentiable with $F'=f$.
  \item The \textbf{Second FTC} (Newton--Leibniz): if $G'=f$ on $(a,b)$,
    then $\int_a^b f = G(b)-G(a)$.
  \item \textbf{Integration by parts} and \textbf{substitution} are the main
    computational tools.
  \item \textbf{Improper integrals} extend the theory to unbounded domains
    or unbounded integrands; convergence is tested by comparison.
\end{itemize}
\end{tcolorbox}

\index{Riemann integral|)}


%%%%%%%%%%%%%%%%%%%%%%%%%%%%%%%%%%%%%%%%%%%%%%%%%%%%%%%%%%%%
%%%%%%%%%%%%%%%%%%%%%%%%%%%%%%%%%%%%%%%%%%%%%%%%%%%%%%%%%%%%
%%
%%  CHAPTER 8 : SEQUENCES AND SERIES OF FUNCTIONS
%%
%%%%%%%%%%%%%%%%%%%%%%%%%%%%%%%%%%%%%%%%%%%%%%%%%%%%%%%%%%%%
%%%%%%%%%%%%%%%%%%%%%%%%%%%%%%%%%%%%%%%%%%%%%%%%%%%%%%%%%%%%

\chapter{Sequences and Series of Functions}\label{ch:seq-series-fn}
\index{sequences of functions|(}
\index{series of functions|(}

\section*{Motivation}
\addcontentsline{toc}{section}{Motivation}

In Analysis~I we studied sequences and series of real numbers.  Now we
consider sequences $(f_n)$ and series $\sum f_n$ whose terms are
\emph{functions}.  The central question is:

\begin{center}
\itshape
Under what conditions can we interchange a limit with an integral,
a derivative, or another limit?
\end{center}

The answer depends on the \emph{mode of convergence}.  Pointwise
convergence alone is too weak to guarantee these interchanges; the
correct notion turns out to be \emph{uniform convergence}.  This chapter
develops the theory systematically, culminating in the Weierstrass
$M$-test, power series, and the Weierstrass approximation theorem.


%% ──────────────────────────────────────────────────────
\section{Pointwise and Uniform Convergence}\label{sec:pointwise-uniform}

\begin{definition}[Pointwise convergence]\label{def:pointwise-conv}
  \index{pointwise convergence}
  Let $(f_n)_{n\ge1}$ be a sequence of functions $f_n\colon D\to\R$.
  We say $(f_n)$ \emph{converges pointwise} to $f\colon D\to\R$ if for
  every $x\in D$,
  \[
    \lim_{n\to\infty}f_n(x) = f(x),
  \]
  i.e., for every $x\in D$ and every $\eps>0$, there exists
  $N=N(x,\eps)\in\N$ such that $n\ge N$ implies $\abs{f_n(x)-f(x)}<\eps$.
\end{definition}

The crucial point in pointwise convergence is that $N$ may depend on $x$.
When we require $N$ to be \emph{independent} of $x$, we obtain a
stronger mode of convergence.

\begin{definition}[Uniform convergence]\label{def:uniform-conv}
  \index{uniform convergence}
  The sequence $(f_n)$ \emph{converges uniformly} to $f$ on $D$ if for
  every $\eps>0$, there exists $N=N(\eps)\in\N$ (independent of~$x$)
  such that
  \[
    n\ge N \implies \sup_{x\in D}\abs{f_n(x)-f(x)} < \eps.
  \]
  Equivalently, $\norm{f_n-f}_\infty\to0$, where
  $\norm{g}_\infty=\sup_{x\in D}\abs{g(x)}$.
\end{definition}

\begin{remark}[Negation of uniform convergence]\label{rem:negation-uniform}
  $(f_n)$ does \emph{not} converge uniformly to $f$ on $D$ if there
  exists $\eps_0>0$ and a sequence $(x_n)$ in $D$ such that
  $\abs{f_n(x_n)-f(x_n)}\ge\eps_0$ for infinitely many $n$.
\end{remark}

\begin{center}
\begin{tikzpicture}
  \begin{axis}[
    axis lines=middle,
    xlabel={$x$}, ylabel={$y$},
    xmin=-0.2, xmax=3.5,
    ymin=-0.5, ymax=2.5,
    width=12cm, height=7cm,
    xtick=\empty, ytick=\empty,
    every axis x label/.style={at={(ticklabel* cs:1)},anchor=west},
    every axis y label/.style={at={(ticklabel* cs:1)},anchor=south},
    title={The $\eps$-tube: uniform convergence means $f_n$ eventually
           lies in the shaded band},
  ]
    % Limit function f
    \addplot[domain=0.1:3.2, samples=100, very thick, black]
      {1.5*exp(-0.3*x)*sin(deg(2*x))+1};
    \node at (axis cs:3.0,0.6) {$f$};
    % eps-tube upper
    \addplot[domain=0.1:3.2, samples=100, thick, blue!50, dashed]
      {1.5*exp(-0.3*x)*sin(deg(2*x))+1+0.3};
    \node at (axis cs:3.35,1.35) {\small$f+\eps$};
    % eps-tube lower
    \addplot[domain=0.1:3.2, samples=100, thick, blue!50, dashed]
      {1.5*exp(-0.3*x)*sin(deg(2*x))+1-0.3};
    \node at (axis cs:3.35,0.65) {\small$f-\eps$};
    % Shaded region
    \addplot[fill=blue!8, draw=none, domain=0.1:3.2, samples=100]
      {1.5*exp(-0.3*x)*sin(deg(2*x))+1+0.3}
      \closedcycle;
    % Re-draw limit function on top
    \addplot[domain=0.1:3.2, samples=100, very thick, black]
      {1.5*exp(-0.3*x)*sin(deg(2*x))+1};
    % f_n inside tube
    \addplot[domain=0.1:3.2, samples=100, thick, red!70!black]
      {1.5*exp(-0.3*x)*sin(deg(2*x))+1 + 0.15*sin(deg(5*x))};
    \node[red!70!black] at (axis cs:2.7,1.65) {$f_n$};
  \end{axis}
\end{tikzpicture}
\end{center}


%% ──────────────────────────────────────────────────────
\section{Uniform Cauchy Criterion}\label{sec:uniform-cauchy}

\begin{theorem}[Uniform Cauchy criterion]\label{thm:uniform-cauchy}
  \index{Cauchy criterion!uniform convergence}
  A sequence $(f_n)$ of bounded functions on $D$ converges uniformly on
  $D$ if and only if for every $\eps>0$ there exists $N\in\N$ such that
  \[
    m,n\ge N \implies \sup_{x\in D}\abs{f_m(x)-f_n(x)} < \eps.
  \]
\end{theorem}

\begin{proof}
  \textbf{($\Rightarrow$).}\;
  If $f_n\to f$ uniformly, then for $\eps>0$ choose $N$ so that
  $n\ge N$ implies $\norm{f_n-f}_\infty<\eps/2$.  For $m,n\ge N$,
  \[
    \abs{f_m(x)-f_n(x)} \le \abs{f_m(x)-f(x)}+\abs{f(x)-f_n(x)}
    < \frac{\eps}{2}+\frac{\eps}{2}=\eps.
  \]

  \textbf{($\Leftarrow$).}\;
  For each $x\in D$, the sequence $(f_n(x))$ is Cauchy in $\R$, hence
  converges; call the limit $f(x)$.  We must show the convergence is
  uniform.

  Given $\eps>0$, choose $N$ from the Cauchy condition.  For $m,n\ge N$
  and any $x\in D$, $\abs{f_m(x)-f_n(x)}<\eps$.  Fixing $n\ge N$ and
  letting $m\to\infty$, we obtain $\abs{f(x)-f_n(x)}\le\eps$ for all
  $x\in D$.  Since this bound is independent of $x$,
  $\norm{f-f_n}_\infty\le\eps$, so $f_n\to f$ uniformly.
\end{proof}


%% ──────────────────────────────────────────────────────
\section{A Fundamental Counterexample}\label{sec:counterexample-xn}

\begin{example}[{$f_n(x)=x^n$ on $[0,1]$}]\label{ex:xn-counterexample}
  \index{counterexample!$x^n$ on $[0,1]$}
  Define $f_n\colon[0,1]\to\R$ by $f_n(x)=x^n$.  Each $f_n$ is
  continuous.

  \emph{Pointwise limit.}\;
  For $x\in[0,1)$, $x^n\to0$; for $x=1$, $x^n=1$.  Hence the pointwise
  limit is
  \[
    f(x) = \begin{cases} 0 & \text{if } 0\le x<1,\\ 1 & \text{if } x=1.
    \end{cases}
  \]
  This limit function is \emph{discontinuous} at $x=1$, even though
  every $f_n$ is continuous.

  \emph{Convergence is not uniform.}\;
  Indeed, $\sup_{x\in[0,1]}\abs{f_n(x)-f(x)}
  =\sup_{x\in[0,1)}x^n=1$ for every $n$ (the supremum is not attained
  but equals~$1$).  So $\norm{f_n-f}_\infty=1\not\to0$.

  Alternatively, take $x_n=(1-1/n)$; then
  $f_n(x_n)=(1-1/n)^n\to1/e\ne0=f(x_n)$.
\end{example}

This example shows that the pointwise limit of continuous functions need
not be continuous.  Uniform convergence is precisely the condition that
prevents such pathologies.


%% ──────────────────────────────────────────────────────
\section{Uniform Limit of Continuous Functions}\label{sec:uniform-cont}

\begin{theorem}[Uniform limit of continuous functions is continuous]\label{thm:unif-limit-cont}
  \index{uniform convergence!preserves continuity}
  Let $(f_n)$ be a sequence of continuous functions $f_n\colon D\to\R$
  that converges uniformly to $f\colon D\to\R$.  Then $f$ is continuous
  on $D$.
\end{theorem}

\begin{proof}
  Fix $x_0\in D$ and let $\eps>0$.  We use the ``$\eps/3$-argument.''

  \emph{Step 1.}\;
  Since $f_n\to f$ uniformly, there exists $N\in\N$ such that
  \[
    \sup_{x\in D}\abs{f_N(x)-f(x)} < \frac{\eps}{3}.
  \]

  \emph{Step 2.}\;
  Since $f_N$ is continuous at $x_0$, there exists $\delta>0$ such that
  \[
    \abs{x-x_0}<\delta \implies \abs{f_N(x)-f_N(x_0)} < \frac{\eps}{3}.
  \]

  \emph{Step 3.}\;
  For $\abs{x-x_0}<\delta$, the triangle inequality gives
  \begin{align*}
    \abs{f(x)-f(x_0)}
    &\le \abs{f(x)-f_N(x)} + \abs{f_N(x)-f_N(x_0)} + \abs{f_N(x_0)-f(x_0)} \\
    &< \frac{\eps}{3}+\frac{\eps}{3}+\frac{\eps}{3} = \eps.
  \end{align*}
  Hence $f$ is continuous at $x_0$.  Since $x_0$ was arbitrary, $f$ is
  continuous on $D$.
\end{proof}

\begin{remark}[Converse fails]\label{rem:converse-fails}
  The converse is false: a sequence of continuous functions may converge
  pointwise to a continuous function without the convergence being
  uniform.  For instance, $f_n(x)=x^n/n$ on $[0,1]$ converges
  pointwise to $f\equiv0$ (which is continuous), but
  $\norm{f_n}_\infty=1/n\to0$---wait, that \emph{is} uniform.  A better
  example: $f_n(x)=nx(1-x)^n$ on $[0,1]$ converges pointwise to~$0$
  (continuous), but $\sup f_n=n\cdot\frac{n}{(n+1)^{n+1}}\cdot
  (n+1)\not\to0$ in general---the point is that uniform convergence is a
  strictly stronger condition, and one must verify it explicitly.
\end{remark}

\begin{example}[Pointwise limit continuous but convergence not uniform]\label{ex:pw-cont-not-unif}
  Let $f_n(x)=\frac{nx}{1+n^2x^2}$ on $[0,1]$.  Then $f_n(x)\to0$
  pointwise for every $x\in[0,1]$, and the limit $f\equiv0$ is
  continuous.  However, $f_n(1/n)=\frac{1}{2}$ for every $n$, so
  $\norm{f_n}_\infty\ge 1/2$ and the convergence is not uniform.
\end{example}


%% ──────────────────────────────────────────────────────
\section{Interchange of Limit and Integral}\label{sec:interchange-int}

\begin{theorem}[Uniform convergence and integration]\label{thm:unif-int}
  \index{uniform convergence!integration interchange}
  Let $(f_n)$ be a sequence of Riemann-integrable functions on $[a,b]$
  converging uniformly to $f$.  Then $f$ is Riemann integrable and
  \[
    \lim_{n\to\infty}\int_a^b f_n(x)\dd x = \int_a^b f(x)\dd x.
  \]
  In short, $\displaystyle\int_a^b\lim_{n\to\infty}f_n
  = \lim_{n\to\infty}\int_a^b f_n$.
\end{theorem}

\begin{proof}
  \emph{Step 1: $f$ is integrable.}\;
  Given $\eps>0$, choose $N$ so that $\norm{f_N-f}_\infty<\eps/(3(b-a))$.
  Since $f_N$ is integrable, there exists a partition $P$ with
  $U(f_N,P)-L(f_N,P)<\eps/3$.

  On any subinterval $[x_{i-1},x_i]$,
  \begin{align*}
    \sup_{[x_{i-1},x_i]}f - \inf_{[x_{i-1},x_i]}f
    &\le \sup(f-f_N) - \inf(f-f_N) + \sup f_N - \inf f_N \\
    &\le 2\norm{f-f_N}_\infty + (M_i^N - m_i^N),
  \end{align*}
  where $M_i^N,m_i^N$ denote the sup and inf of $f_N$ on the subinterval.
  Therefore
  \begin{align*}
    U(f,P)-L(f,P)
    &\le 2\norm{f-f_N}_\infty(b-a) + U(f_N,P)-L(f_N,P) \\
    &< 2\cdot\frac{\eps}{3(b-a)}\cdot(b-a) + \frac{\eps}{3} = \eps.
  \end{align*}
  By the integrability criterion, $f\in\mathcal{R}([a,b])$.

  \emph{Step 2: convergence of integrals.}\;
  For $n\ge N$,
  \[
    \abs*{\int_a^b f_n\dd x - \int_a^b f\dd x}
    = \abs*{\int_a^b(f_n-f)\dd x}
    \le \int_a^b\abs{f_n-f}\dd x
    \le \norm{f_n-f}_\infty(b-a).
  \]
  Since $\norm{f_n-f}_\infty\to0$, the right-hand side tends to zero.
\end{proof}

\begin{example}[Failed interchange without uniform convergence]\label{ex:failed-int-interchange}
  \index{counterexample!limit-integral interchange}
  Let $f_n\colon[0,1]\to\R$ be defined by
  \[
    f_n(x) = \begin{cases}
      n^2 x & \text{if } 0\le x\le 1/n,\\
      2n - n^2 x & \text{if } 1/n < x\le 2/n,\\
      0 & \text{if } 2/n < x\le 1.
    \end{cases}
  \]
  Then $f_n(x)\to0$ pointwise for every $x\in[0,1]$, but
  $\int_0^1 f_n\dd x = 1$ for every $n$.  The interchange fails because
  the convergence is not uniform (in fact $\sup f_n=n\to\infty$).
\end{example}


%% ──────────────────────────────────────────────────────
\section{Interchange of Limit and Derivative}\label{sec:interchange-deriv}

\begin{theorem}[Uniform convergence and differentiation]\label{thm:unif-diff}
  \index{uniform convergence!differentiation interchange}
  Let $(f_n)$ be a sequence of $C^1$ functions on $[a,b]$.  Suppose that:
  \begin{enumerate}[label=(\roman*)]
    \item There exists $x_0\in[a,b]$ such that $(f_n(x_0))$ converges.
    \item The sequence of derivatives $(f_n')$ converges uniformly on
      $[a,b]$ to some function $g$.
  \end{enumerate}
  Then $(f_n)$ converges uniformly on $[a,b]$ to a $C^1$ function $f$,
  and $f'=g$.  That is,
  \[
    \left(\lim_{n\to\infty}f_n\right)' = \lim_{n\to\infty}f_n'.
  \]
\end{theorem}

\begin{proof}
  \emph{Step 1: uniform convergence of $(f_n)$.}\;
  By the Fundamental Theorem of Calculus,
  \[
    f_n(x) = f_n(x_0) + \int_{x_0}^x f_n'(t)\dd t.
  \]
  Define $f(x) = \ell + \int_{x_0}^x g(t)\dd t$, where
  $\ell=\lim_{n\to\infty}f_n(x_0)$.

  For any $x\in[a,b]$,
  \begin{align*}
    \abs{f_n(x)-f(x)}
    &\le \abs{f_n(x_0)-\ell} + \abs*{\int_{x_0}^x\bigl(f_n'(t)-g(t)\bigr)\dd t}\\
    &\le \abs{f_n(x_0)-\ell} + \norm{f_n'-g}_\infty\cdot(b-a).
  \end{align*}
  Both terms on the right tend to zero (the first by hypothesis~(i), the
  second by hypothesis~(ii)), independently of~$x$.  Hence $f_n\to f$
  uniformly.

  \emph{Step 2: $f$ is $C^1$ and $f'=g$.}\;
  Since $g$ is the uniform limit of continuous functions $f_n'$, it is
  continuous by Theorem~\ref{thm:unif-limit-cont}.  By the First
  Fundamental Theorem of Calculus,
  $\displaystyle f(x)=\ell+\int_{x_0}^x g(t)\dd t$ is differentiable with
  $f'(x)=g(x)$.
\end{proof}

\begin{remark}[Hypotheses are sharp]\label{rem:deriv-interchange-sharp}
  The condition that $(f_n')$ converges \emph{uniformly} cannot be
  weakened to pointwise.  Mere pointwise convergence of $(f_n)$ and
  $(f_n')$ does not ensure $(\lim f_n)'=\lim f_n'$.  For instance,
  $f_n(x)=\frac{\sin(nx)}{\sqrt{n}}$ converges uniformly to~$0$ on
  $\R$, but $f_n'(x)=\sqrt{n}\cos(nx)$ does not converge at all.
\end{remark}


%% ──────────────────────────────────────────────────────
\section{Series of Functions}\label{sec:series-fn}

\begin{definition}[Pointwise, uniform, and normal convergence of series]\label{def:series-modes}
  \index{series of functions!modes of convergence}
  Let $(u_n)$ be a sequence of functions $u_n\colon D\to\R$.  The series
  $\sum_{n\ge0}u_n$ is said to converge:
  \begin{itemize}
    \item \textbf{Pointwise} on $D$ if for every $x\in D$, the numerical
      series $\sum u_n(x)$ converges.
    \item \textbf{Uniformly} on $D$ if the sequence of partial sums
      $S_N=\sum_{n=0}^N u_n$ converges uniformly on $D$.
    \item \textbf{Normally}\index{normal convergence} on $D$ if
      $\sum_{n\ge0}\norm{u_n}_\infty<+\infty$,
      where $\norm{u_n}_\infty=\sup_{x\in D}\abs{u_n(x)}$.
  \end{itemize}
\end{definition}

\begin{proposition}[Normal convergence implies uniform convergence]\label{prop:normal-implies-unif}
  If $\sum u_n$ converges normally on $D$, then it converges uniformly
  (and hence pointwise) on $D$.
\end{proposition}

\begin{proof}
  For $N<M$,
  $\norm{S_M-S_N}_\infty = \norm{\sum_{n=N+1}^M u_n}_\infty
  \le \sum_{n=N+1}^M \norm{u_n}_\infty$.  Since $\sum\norm{u_n}_\infty$
  converges, its tail tends to zero, so $(S_N)$ is uniformly Cauchy,
  hence uniformly convergent.
\end{proof}


%% ──────────────────────────────────────────────────────
\section{The Weierstrass \texorpdfstring{$M$}{M}-Test}\label{sec:M-test}

\begin{theorem}[Weierstrass $M$-test]\label{thm:M-test}
  \index{Weierstrass $M$-test}
  Let $(u_n)$ be a sequence of functions on $D$ and suppose there exist
  constants $M_n\ge0$ with $\abs{u_n(x)}\le M_n$ for all $x\in D$ and
  $\sum_{n\ge0}M_n<+\infty$.  Then $\sum u_n$ converges uniformly
  (and absolutely) on $D$.
\end{theorem}

\begin{proof}
  We have $\norm{u_n}_\infty\le M_n$, so
  $\sum\norm{u_n}_\infty\le\sum M_n<+\infty$.  Hence $\sum u_n$
  converges normally, and by Proposition~\ref{prop:normal-implies-unif},
  it converges uniformly.

  Absolute convergence at each $x$ follows from
  $\sum\abs{u_n(x)}\le\sum M_n<+\infty$.
\end{proof}

\begin{example}[Application of the $M$-test]\label{ex:M-test-example}
  The series $\displaystyle\sum_{n=1}^{\infty}\frac{\sin(nx)}{n^2}$
  converges uniformly on $\R$.  Indeed,
  $\abs{\frac{\sin(nx)}{n^2}}\le\frac{1}{n^2}$ and
  $\sum\frac{1}{n^2}<+\infty$.
\end{example}


%% ──────────────────────────────────────────────────────
\section{Interchange Theorems for Series}\label{sec:interchange-series}

The interchange theorems for sequences translate directly into theorems
for series via partial sums.

\begin{corollary}[Term-by-term integration]\label{cor:termwise-int}
  \index{series of functions!term-by-term integration}
  If $\sum u_n$ converges uniformly on $[a,b]$ and each $u_n$ is Riemann
  integrable, then $\sum u_n$ is integrable and
  \[
    \int_a^b\sum_{n=0}^{\infty}u_n(x)\dd x
    = \sum_{n=0}^{\infty}\int_a^b u_n(x)\dd x.
  \]
\end{corollary}

\begin{proof}
  Apply Theorem~\ref{thm:unif-int} to the partial sums
  $S_N=\sum_{n=0}^N u_n$.
\end{proof}

\begin{corollary}[Term-by-term differentiation]\label{cor:termwise-diff}
  \index{series of functions!term-by-term differentiation}
  Suppose each $u_n$ is $C^1$ on $[a,b]$, $\sum u_n(x_0)$ converges for
  some $x_0\in[a,b]$, and $\sum u_n'$ converges uniformly on $[a,b]$.
  Then $\sum u_n$ converges uniformly to a $C^1$ function $S$ and
  \[
    S'(x) = \sum_{n=0}^{\infty}u_n'(x).
  \]
\end{corollary}

\begin{proof}
  Apply Theorem~\ref{thm:unif-diff} to the partial sums.
\end{proof}


%% ──────────────────────────────────────────────────────
\section{Power Series}\label{sec:power-series}
\index{power series|(}

\begin{definition}[Power series]\label{def:power-series-def}
  A \emph{power series} centered at $a\in\R$ is a series of the form
  \[
    \sum_{n=0}^{\infty}c_n(x-a)^n,
  \]
  where $(c_n)_{n\ge0}$ is a sequence of real coefficients.
\end{definition}

\begin{theorem}[Radius of convergence]\label{thm:radius-conv}
  \index{power series!radius of convergence}
  For every power series $\sum c_n(x-a)^n$ there exists a unique
  $R\in[0,+\infty]$ (the \emph{radius of convergence}) such that the
  series converges absolutely for $\abs{x-a}<R$ and diverges for
  $\abs{x-a}>R$.  Moreover,
  \[
    \frac{1}{R} = \limsup_{n\to\infty}\abs{c_n}^{1/n}
    \qquad\text{(Cauchy--Hadamard formula).}
  \]
\end{theorem}

\begin{theorem}[Uniform convergence of power series on compact subsets]\label{thm:power-unif-compact}
  \index{power series!uniform convergence}
  Let $\sum c_n(x-a)^n$ have radius of convergence $R>0$.  Then for
  every $0<r<R$, the series converges uniformly (in fact normally) on
  $[a-r,a+r]$.
\end{theorem}

\begin{proof}
  For $\abs{x-a}\le r<R$, $\abs{c_n(x-a)^n}\le\abs{c_n}r^n$.  Since
  $r<R$, the series $\sum\abs{c_n}r^n$ converges (by definition of the
  radius of convergence).  The Weierstrass $M$-test with
  $M_n=\abs{c_n}r^n$ gives uniform convergence on $[a-r,a+r]$.
\end{proof}

\begin{theorem}[Term-by-term differentiation and integration of power series]\label{thm:power-diff-int}
  \index{power series!term-by-term differentiation}
  \index{power series!term-by-term integration}
  Let $f(x)=\sum_{n=0}^{\infty}c_n(x-a)^n$ with radius of convergence
  $R>0$.  Then:
  \begin{enumerate}[label=(\alph*)]
    \item $f$ is infinitely differentiable on $(a-R,a+R)$ and
      \[
        f'(x) = \sum_{n=1}^{\infty}n\,c_n(x-a)^{n-1},
      \]
      with the derived series having the same radius of convergence $R$.
    \item For any $x\in(a-R,a+R)$,
      \[
        \int_a^x f(t)\dd t
        = \sum_{n=0}^{\infty}\frac{c_n}{n+1}(x-a)^{n+1},
      \]
      with the integrated series having the same radius of convergence $R$.
  \end{enumerate}
\end{theorem}

\begin{proof}
  \textbf{(a).}\;
  Let $g(x)=\sum_{n=1}^{\infty}nc_n(x-a)^{n-1}$.  We first show that
  the derived series has radius of convergence $R$.  Since
  $\limsup_{n\to\infty}\abs{nc_n}^{1/n}
  =\limsup\abs{c_n}^{1/n}\cdot\lim n^{1/n}
  =\limsup\abs{c_n}^{1/n}\cdot1=1/R$, the Cauchy--Hadamard formula
  gives the same radius.

  Now fix $x_0\in(a-R,a+R)$ and choose $r$ with $\abs{x_0-a}<r<R$.
  On $[a-r,a+r]$, the series $\sum nc_n(x-a)^{n-1}$ converges uniformly
  (by the $M$-test with $M_n=n\abs{c_n}r^{n-1}$, since $\sum M_n$
  converges).  Also, $\sum c_n(x_0-a)^n$ converges.  By
  Corollary~\ref{cor:termwise-diff}, $f$ is differentiable and
  $f'(x)=g(x)$ on $(a-r,a+r)$.  Since $r$ can be taken arbitrarily
  close to $R$, this holds on all of $(a-R,a+R)$.  Iterating gives
  infinite differentiability.

  \textbf{(b).}\;
  The integrated series has the same radius by Cauchy--Hadamard (since
  $\abs{c_n/(n+1)}^{1/n}\to\abs{c_n}^{1/n}$ in the limsup).  On any
  compact sub-interval of $(a-R,a+R)$ the original series converges
  uniformly, so term-by-term integration is justified by
  Corollary~\ref{cor:termwise-int}.
\end{proof}


%% ──────────────────────────────────────────────────────
\section{Abel's Theorem}\label{sec:abel}

\begin{theorem}[Abel's theorem]\label{thm:abel}
  \index{Abel's theorem}
  Let $\sum_{n=0}^{\infty}c_n x^n$ have radius of convergence $R>0$
  (which may be $+\infty$).  If the series $\sum_{n=0}^{\infty}c_n R^n$
  converges (i.e.\ the power series converges at the right endpoint $x=R$),
  then
  \[
    \lim_{x\to R^-}\sum_{n=0}^{\infty}c_n x^n
    = \sum_{n=0}^{\infty}c_n R^n.
  \]
  That is, the function $f(x)=\sum c_n x^n$ is left-continuous at $x=R$.
  An analogous statement holds at $x=-R$.
\end{theorem}

Abel's theorem is a delicate result whose proof uses Abel summation
(summation by parts).  It has important applications, for instance in
deriving
$\ln 2 = 1-\frac{1}{2}+\frac{1}{3}-\frac{1}{4}+\cdots$ from the power
series $\ln(1+x)=\sum_{n=1}^{\infty}\frac{(-1)^{n+1}}{n}x^n$
($R=1$, converges at $x=1$ by the alternating series test).

\index{power series|)}


%% ──────────────────────────────────────────────────────
\section{The Weierstrass Approximation Theorem}\label{sec:weierstrass-approx}
\index{Weierstrass approximation theorem}

\begin{theorem}[Weierstrass approximation theorem]\label{thm:weierstrass-approx}
  Every continuous function $f\colon[a,b]\to\R$ can be uniformly
  approximated by polynomials.  That is, for every $\eps>0$ there exists
  a polynomial $p$ such that
  \[
    \sup_{x\in[a,b]}\abs{f(x)-p(x)} < \eps.
  \]
  Equivalently, the polynomials are dense in $(C([a,b]),\norm{\cdot}_\infty)$.
\end{theorem}

\noindent\textbf{Proof idea via Bernstein polynomials.}\;
\index{Bernstein polynomials}
Without loss of generality take $[a,b]=[0,1]$.  For $f\in C([0,1])$
and $n\ge1$, the $n$-th \emph{Bernstein polynomial} of $f$ is
\[
  B_n(f)(x) = \sum_{k=0}^n f\!\left(\frac{k}{n}\right)\binom{n}{k}
  x^k(1-x)^{n-k}.
\]
This is a polynomial of degree at most $n$ that samples $f$ at the
points $k/n$ and takes a weighted average using the binomial
distribution.

The key properties used in the proof are:
\begin{enumerate}[label=(\roman*)]
  \item $B_n(1)(x)=1$ (the binomial theorem),
  \item $B_n(t)(x)=x$,
  \item $B_n(t^2)(x)=x^2+\frac{x(1-x)}{n}$.
\end{enumerate}
Using the uniform continuity of $f$ on $[0,1]$ and properties (i)--(iii)
to control the variance, one shows that
$\norm{B_n(f)-f}_\infty\to0$ as $n\to\infty$.

\begin{center}
\begin{tikzpicture}
  \begin{axis}[
    axis lines=middle,
    xlabel={$x$}, ylabel={$y$},
    xmin=-0.1, xmax=1.15,
    ymin=-0.2, ymax=1.8,
    width=11cm, height=7cm,
    xtick={0,0.5,1},
    ytick=\empty,
    every axis x label/.style={at={(ticklabel* cs:1)},anchor=west},
    every axis y label/.style={at={(ticklabel* cs:1)},anchor=south},
    title={Bernstein polynomial approximation of $f(x)=\abs{2x-1}$},
    legend pos=north west,
    legend style={font=\small},
  ]
    % f(x)=|2x-1|
    \addplot[domain=0:0.5, samples=50, very thick, black] {1-2*x};
    \addplot[domain=0.5:1, samples=50, very thick, black] {2*x-1};
    \addlegendentry{$f(x)=|2x-1|$}
    % B_4
    \addplot[domain=0:1, samples=100, thick, blue!70!black, dashed]
      {abs(0-1)*1*(1-x)^4 + abs(2*0.25-1)*4*x*(1-x)^3
       + abs(2*0.5-1)*6*x^2*(1-x)^2 + abs(2*0.75-1)*4*x^3*(1-x)
       + abs(2*1-1)*x^4};
    \addlegendentry{$B_4(f)$}
    % B_10 (approximate)
    \addplot[domain=0:1, samples=200, thick, red!70!black, densely dotted]
      { (1)*(1-x)^10
        + (0.8)*10*x*(1-x)^9
        + (0.6)*45*x^2*(1-x)^8
        + (0.4)*120*x^3*(1-x)^7
        + (0.2)*210*x^4*(1-x)^6
        + (0.0)*252*x^5*(1-x)^5
        + (0.2)*210*x^6*(1-x)^4
        + (0.4)*120*x^7*(1-x)^3
        + (0.6)*45*x^8*(1-x)^2
        + (0.8)*10*x^9*(1-x)
        + (1)*x^10};
    \addlegendentry{$B_{10}(f)$}
  \end{axis}
\end{tikzpicture}
\end{center}


%% ──────────────────────────────────────────────────────
\section{Counterexamples for Failed Interchanges}\label{sec:failed-interchanges}

We collect several instructive counterexamples showing that the
hypotheses of the interchange theorems cannot be relaxed.

\begin{example}[Pointwise convergence does not preserve continuity]\label{ex:counter-cont}
  As seen in Example~\ref{ex:xn-counterexample}, $f_n(x)=x^n$ on $[0,1]$
  converges pointwise to a discontinuous function.
\end{example}

\begin{example}[Pointwise convergence and integration]\label{ex:counter-int}
  \index{counterexample!limit-integral interchange}
  Define $f_n\colon[0,1]\to\R$ by $f_n(x)=n\cdot\mathbf{1}_{(0,1/n]}(x)$.
  Then $f_n\to0$ pointwise, but $\int_0^1 f_n\dd x=1$ for all $n$.
\end{example}

\begin{example}[Pointwise convergence and differentiation]\label{ex:counter-diff}
  \index{counterexample!limit-derivative interchange}
  Let $f_n(x)=\frac{\sin(nx)}{n}$ on $\R$.  Then $f_n\to0$ uniformly,
  but $f_n'(x)=\cos(nx)$ does not converge at all (e.g.\ at $x=0$,
  $f_n'(0)=1$ for every $n$, while at $x=\pi$, $f_n'(\pi)=\cos(n\pi)=(-1)^n$).
  Note that $(f_n')$ does \emph{not} converge uniformly, so the
  hypotheses of Theorem~\ref{thm:unif-diff} are not met.
\end{example}

\begin{example}[Differentiation of a pointwise but non-uniformly convergent series]\label{ex:counter-series-diff}
  The series $\sum_{n=1}^{\infty}\frac{x}{n(1+nx^2)}$ converges
  pointwise on $\R$ (compare with $1/n^2$ for $x\ne0$).  Each term is
  $C^1$, but the series of derivatives
  $\sum\frac{1-nx^2}{n(1+nx^2)^2}$ diverges at $x=0$.  This shows that
  without uniform convergence of the derived series, term-by-term
  differentiation fails.
\end{example}


%% ──────────────────────────────────────────────────────
\section{Exercises}\label{sec:exer-ch8}

\begin{exercise}[Uniform convergence on sub-domains]\label{exo:exer-unif-sub}
  Show that $f_n(x)=x^n$ converges uniformly on $[0,a]$ for every
  $0<a<1$, but not uniformly on $[0,1]$.
\end{exercise}

\begin{exercise}[Sup-norm computation]\label{exo:exer-sup-norm}
  Let $f_n(x)=\frac{x}{1+nx^2}$ on $[0,+\infty)$.  Compute
  $\norm{f_n}_\infty$ and determine whether $(f_n)$ converges uniformly
  on $[0,+\infty)$.
\end{exercise}

\begin{exercise}[Uniform Cauchy criterion application]\label{exo:exer-cauchy-apply}
  Show that $f_n(x)=\sum_{k=0}^n\frac{x^k}{k!}$ converges uniformly on
  every interval $[-R,R]$, $R>0$, by verifying the uniform Cauchy
  criterion.
\end{exercise}

\begin{exercise}[Uniform limit preserves boundedness]\label{exo:exer-bdd}
  Prove that if each $f_n$ is bounded on $D$ and $f_n\to f$ uniformly
  on $D$, then $f$ is bounded on $D$.
\end{exercise}

\begin{exercise}[Weierstrass $M$-test applications]\label{exo:exer-M-test}
  Show that each series converges uniformly on the given domain:
  \begin{enumerate}[label=(\alph*)]
    \item $\displaystyle\sum_{n=1}^{\infty}\frac{\cos(nx)}{n^3}$ on $\R$
    \item $\displaystyle\sum_{n=0}^{\infty}\frac{x^n}{n!}$ on $[-R,R]$
      for any $R>0$
    \item $\displaystyle\sum_{n=1}^{\infty}\frac{(-1)^n}{n+x^2}$ on $\R$
      (use Leibniz and then verify uniformity)
  \end{enumerate}
\end{exercise}

\begin{exercise}[Continuity of a series]\label{exo:exer-series-cont}
  Show that $\displaystyle f(x)=\sum_{n=1}^{\infty}\frac{1}{n^2+x^2}$
  is continuous on $\R$.
\end{exercise}

\begin{exercise}[Term-by-term integration]\label{exo:exer-termwise-int}
  Let $f(x)=\sum_{n=0}^{\infty}(-1)^n x^{2n}=\frac{1}{1+x^2}$ for
  $\abs{x}<1$.  By integrating term by term, derive the Leibniz formula
  $\displaystyle\frac{\pi}{4}=\sum_{n=0}^{\infty}\frac{(-1)^n}{2n+1}$.
  Justify the interchange at $x=1$ using Abel's theorem.
\end{exercise}

\begin{exercise}[Interchange failure for derivatives]\label{exo:exer-diff-fail}
  Let $f_n(x)=\frac{1}{n}\sin(n^2 x)$.  Show that $f_n\to0$ uniformly
  on $\R$, but $(f_n')$ does not converge at any $x$.  Why does this not
  contradict Theorem~\ref{thm:unif-diff}?
\end{exercise}

\begin{exercise}[Power series and radius]\label{exo:exer-radius}
  Find the radius of convergence of each power series:
  \begin{enumerate}[label=(\alph*)]
    \item $\displaystyle\sum_{n=0}^{\infty}\frac{n!}{n^n}x^n$
    \item $\displaystyle\sum_{n=1}^{\infty}\frac{x^n}{n\cdot 2^n}$
    \item $\displaystyle\sum_{n=0}^{\infty}\frac{(2n)!}{(n!)^2}x^n$
  \end{enumerate}
\end{exercise}

\begin{exercise}[Abel's theorem application]\label{exo:exer-abel}
  Using the power series $\ln(1+x)=\sum_{n=1}^{\infty}\frac{(-1)^{n+1}}{n}x^n$
  ($\abs{x}<1$) and Abel's theorem, prove that
  $\ln 2=1-\frac{1}{2}+\frac{1}{3}-\frac{1}{4}+\cdots$.
\end{exercise}

\begin{exercise}[Uniform but not normal convergence]\label{exo:exer-unif-not-normal}
  Show that $\sum_{n=1}^{\infty}\frac{(-1)^n}{n+x^2}$ converges
  uniformly on $[0,+\infty)$ but not normally.
  \emph{Hint:} the sup-norm of the $n$-th term is $1/n$.
\end{exercise}

\begin{exercise}[Dini's theorem]\label{exo:exer-dini}
  \index{Dini's theorem}
  Let $(f_n)$ be a sequence of continuous functions on a compact set
  $K$ converging pointwise to a continuous function $f$.  If $f_n$
  is \emph{monotone} (i.e.\ $f_1\ge f_2\ge\cdots$ or
  $f_1\le f_2\le\cdots$), prove that the convergence is uniform.
  \emph{Hint:} consider $g_n=f_n-f$ (or $f-f_n$); use compactness to
  extract a finite subcover from the open cover
  $\{x:\abs{g_n(x)}<\eps\}$.
\end{exercise}

\bigskip
\begin{tcolorbox}[colback=gray!5, colframe=gray!60, title={\textbf{Chapter 8 --- Summary}}]
\begin{itemize}[leftmargin=*]
  \item \textbf{Pointwise convergence}: $f_n(x)\to f(x)$ for each $x$;
    the rate may depend on $x$.
  \item \textbf{Uniform convergence}: $\norm{f_n-f}_\infty\to0$; the
    rate is independent of $x$.
  \item The uniform limit of continuous functions is \textbf{continuous}
    ($\eps/3$-argument).
  \item Under uniform convergence, one may \textbf{interchange limit and
    integral}.
  \item To interchange \textbf{limit and derivative}, one needs uniform
    convergence of the \emph{derivatives} (plus pointwise convergence of
    the functions at one point).
  \item The \textbf{Weierstrass $M$-test} gives a practical criterion for
    uniform (in fact normal) convergence of series.
  \item \textbf{Power series} converge uniformly on compact subsets of
    their interval of convergence and can be differentiated and
    integrated term by term.
  \item \textbf{Abel's theorem} ensures continuity of a power series at
    an endpoint where it converges.
  \item The \textbf{Weierstrass approximation theorem}: every continuous
    function on $[a,b]$ is the uniform limit of polynomials.
  \item \textbf{Counterexamples} show that pointwise convergence alone
    does \emph{not} justify interchanging limits with continuity,
    integration, or differentiation.
\end{itemize}
\end{tcolorbox}

\index{sequences of functions|)}
\index{series of functions|)}
%%%%%%%%%%%%%%%%%%%%%%%%%%%%%%%%%%%%%%%%%%%%%%%%%%%%%%%%%%%%
%% Real Analysis II — Chapters 9, 10, 11
%% Part 2 (Final): Multivariable Calculus, Multiple Integrals,
%%                  Curves and Line Integrals
%%%%%%%%%%%%%%%%%%%%%%%%%%%%%%%%%%%%%%%%%%%%%%%%%%%%%%%%%%%%

%% ══════════════════════════════════════════════════════════
%%  CHAPTER 9
%% ══════════════════════════════════════════════════════════

\chapter{Introduction to Multivariable Differential Calculus}
\label{ch:multivariable-calculus}
\index{multivariable calculus}

\section*{Introduction}

Throughout the preceding chapters we have studied functions of a single real
variable.  Yet the phenomena of physics, economics, and geometry are almost
never governed by a single parameter.  The temperature at a point in a room
depends on three spatial coordinates and on time; the profit of a firm depends
on quantities produced, prices, and advertising expenditures; the shape of a
surface is described by a function $f(x,y)$.

This chapter extends differential calculus to functions
$f\colon \mathbb{R}^n \to \mathbb{R}^m$.  The central idea is that a
\emph{differentiable} function is one that admits a good \emph{linear}
approximation near each point.  Making this precise requires us first to
discuss the topology of~$\mathbb{R}^n$.

%% ──────────────────────────────────────────────────────────
\section{Topology of \texorpdfstring{$\mathbb{R}^n$}{Rn}}
\label{sec:topology-Rn}
\index{topology!of Rn@of $\mathbb{R}^n$}

\subsection{Norms on \texorpdfstring{$\mathbb{R}^n$}{Rn}}
\label{subsec:norms-Rn}

\begin{definition}{Norm on $\mathbb{R}^n$}{norm-Rn}
A \emph{norm}\index{norm} on $\mathbb{R}^n$ is a function
$\|\cdot\|\colon \mathbb{R}^n \to [0,+\infty)$ satisfying, for all
$x,y\in\mathbb{R}^n$ and $\lambda\in\mathbb{R}$:
\begin{enumerate}[label=(\roman*)]
  \item \textbf{Positive definiteness:} $\|x\|=0 \iff x=0$.
  \item \textbf{Absolute homogeneity:} $\|\lambda x\|=|\lambda|\,\|x\|$.
  \item \textbf{Triangle inequality:} $\|x+y\|\le \|x\|+\|y\|$.
\end{enumerate}
\end{definition}

\begin{example}[Common norms on $\mathbb{R}^n$]
\label{ex:common-norms}
Let $x=(x_1,\dots,x_n)\in\mathbb{R}^n$.
\begin{enumerate}[label=(\alph*)]
  \item \textbf{Euclidean norm}\index{norm!Euclidean}:
    $\|x\|_2 = \bigl(\sum_{i=1}^n x_i^2\bigr)^{1/2}$.
  \item \textbf{1-norm}\index{norm!1-norm}:
    $\|x\|_1 = \sum_{i=1}^n |x_i|$.
  \item \textbf{Sup-norm}\index{norm!sup-norm}:
    $\|x\|_\infty = \max_{1\le i\le n}|x_i|$.
\end{enumerate}
\end{example}

\begin{proposition}[Relationships between norms]
\label{prop:norm-relationships}
For every $x\in\mathbb{R}^n$,
\[
  \|x\|_\infty \;\le\; \|x\|_2 \;\le\; \|x\|_1
  \;\le\; n\,\|x\|_\infty.
\]
\end{proposition}

\begin{proof}
For the first inequality, $\|x\|_\infty = \max_i |x_i|$ and
$\|x\|_2 = (\sum_i x_i^2)^{1/2} \ge (\max_i x_i^2)^{1/2} = \|x\|_\infty$.
For the second, by Cauchy--Schwarz,
$\|x\|_1 = \sum_i |x_i| \le \sqrt{n}\,(\sum_i x_i^2)^{1/2} = \sqrt{n}\,\|x\|_2$;
the sharper bound $\|x\|_2 \le \|x\|_1$ follows from
$(\sum_i |x_i|)^2 = \sum_i x_i^2 + 2\sum_{i<j}|x_i||x_j| \ge \sum_i x_i^2$.
Finally, $\|x\|_1 = \sum_i |x_i| \le n\max_i|x_i| = n\|x\|_\infty$.
\end{proof}

\begin{theorem}[Equivalence of norms in finite dimension]
\label{thm:equiv-norms}
\index{equivalence of norms}
All norms on $\mathbb{R}^n$ are equivalent: for any two norms $\|\cdot\|_a$
and $\|\cdot\|_b$ on $\mathbb{R}^n$, there exist constants $0<c\le C$ such
that
\[
  c\,\|x\|_a \;\le\; \|x\|_b \;\le\; C\,\|x\|_a
  \qquad \text{for all } x\in\mathbb{R}^n.
\]
\end{theorem}

\begin{remark}
The proof uses compactness of the unit sphere in $(\mathbb{R}^n,\|\cdot\|_2)$
and continuity of any norm with respect to the Euclidean norm.  We refer to
Rudin~\cite{rudin} for details.  The conclusion is that in $\mathbb{R}^n$ the
choice of norm does not affect notions of convergence, continuity, or openness.
\end{remark}

\subsection{Open sets, closed sets, and balls}
\label{subsec:open-closed-Rn}

\begin{definition}[Open ball]\label{def:open-ball-Rn}
\index{open ball!in Rn@in $\mathbb{R}^n$}
Given a norm $\|\cdot\|$ on $\mathbb{R}^n$, the \emph{open ball} of centre
$a\in\mathbb{R}^n$ and radius $r>0$ is
\[
  B(a,r) = \{x\in\mathbb{R}^n : \|x-a\| < r\}.
\]
\end{definition}

\begin{definition}{Open and closed sets in $\mathbb{R}^n$}{open-closed-Rn}
\index{open set!in Rn@in $\mathbb{R}^n$}
\index{closed set!in Rn@in $\mathbb{R}^n$}
A set $U\subseteq\mathbb{R}^n$ is \emph{open} if for every $a\in U$ there
exists $r>0$ such that $B(a,r)\subseteq U$.  A set $F$ is \emph{closed}
if its complement $\mathbb{R}^n\setminus F$ is open.
\end{definition}

\begin{center}
\begin{tikzpicture}[scale=1.8]
  % Unit balls in R^2 for different norms
  % sup-norm ball (square)
  \draw[thick, blue!70!black, fill=blue!8]
    (-1,-1) rectangle (1,1);
  \node[blue!70!black] at (0,-1.25) {$\|x\|_\infty \le 1$};

  \begin{scope}[xshift=3.2cm]
  % Euclidean ball (circle)
  \draw[thick, red!70!black, fill=red!8]
    (0,0) circle (1);
  \node[red!70!black] at (0,-1.25) {$\|x\|_2 \le 1$};
  \end{scope}

  \begin{scope}[xshift=6.4cm]
  % 1-norm ball (diamond)
  \draw[thick, green!50!black, fill=green!8]
    (0,1) -- (1,0) -- (0,-1) -- (-1,0) -- cycle;
  \node[green!50!black] at (0,-1.25) {$\|x\|_1 \le 1$};
  \end{scope}
\end{tikzpicture}
\captionof{figure}{Unit balls in $\mathbb{R}^2$ for the sup-norm, Euclidean
norm, and 1-norm.}
\label{fig:unit-balls}
\end{center}

\subsection{Sequences in \texorpdfstring{$\mathbb{R}^n$}{Rn}}
\label{subsec:sequences-Rn}

\begin{proposition}[Component-wise convergence]
\label{prop:component-convergence}
\index{convergence!component-wise}
Let $(x^{(k)})_{k\ge 1}$ be a sequence in $\mathbb{R}^n$ with
$x^{(k)}=(x_1^{(k)},\dots,x_n^{(k)})$, and let $a=(a_1,\dots,a_n)$.  Then
\[
  x^{(k)}\to a \text{ in } \mathbb{R}^n
  \quad\iff\quad
  x_i^{(k)}\to a_i \text{ in } \mathbb{R},\; \forall\, i=1,\dots,n.
\]
\end{proposition}

\begin{proof}
$(\Rightarrow)$\; For each $i$, $|x_i^{(k)}-a_i| \le \|x^{(k)}-a\|_2
\to 0$.

$(\Leftarrow)$\; $\|x^{(k)}-a\|_\infty = \max_i |x_i^{(k)}-a_i| \to 0$,
and $\|x^{(k)}-a\|_2 \le \sqrt{n}\,\|x^{(k)}-a\|_\infty \to 0$.
\end{proof}

%% ──────────────────────────────────────────────────────────
\section{Limits and Continuity in \texorpdfstring{$\mathbb{R}^n$}{Rn}}
\label{sec:limits-continuity-Rn}
\index{limit!in Rn@in $\mathbb{R}^n$}
\index{continuity!in Rn@in $\mathbb{R}^n$}

\begin{definition}[Limit of a function of several variables]\label{def:limit-Rn}
Let $U\subseteq\mathbb{R}^n$ be open, $a\in U$, and
$f\colon U\setminus\{a\}\to\mathbb{R}^m$.  We say
$\lim_{x\to a} f(x)=\ell\in\mathbb{R}^m$ if for every $\varepsilon>0$
there exists $\delta>0$ such that
\[
  0<\|x-a\|<\delta \;\Longrightarrow\; \|f(x)-\ell\|<\varepsilon.
\]
\end{definition}

\begin{definition}[Continuity]\label{def:continuity-Rn}
A function $f\colon U\to\mathbb{R}^m$ is \emph{continuous at $a\in U$} if
$\lim_{x\to a}f(x)=f(a)$.  It is \emph{continuous on $U$} if it is
continuous at every point of~$U$.
\end{definition}

\begin{example}[Iterated limits need not equal the double limit]
\label{ex:iterated-limits}
\index{iterated limits}
Define $f\colon\mathbb{R}^2\setminus\{(0,0)\}\to\mathbb{R}$ by
\[
  f(x,y)=\frac{xy}{x^2+y^2}.
\]
For every fixed $y\ne 0$, $\lim_{x\to 0}f(x,y)=0$; likewise
$\lim_{y\to 0}f(x,y)=0$ for fixed $x\ne 0$.  Both iterated limits equal~$0$.
However, along the line $y=x$, $f(x,x)=\tfrac{x^2}{2x^2}=\tfrac12$, while
along $y=-x$, $f(x,-x)=-\tfrac12$.  Therefore the double limit
$\lim_{(x,y)\to(0,0)}f(x,y)$ does not exist.
\end{example}

\begin{theorem}[Extreme value theorem on compact sets]
\label{thm:extreme-value-Rn}
\index{extreme value theorem!in Rn@in $\mathbb{R}^n$}
Let $K\subseteq\mathbb{R}^n$ be compact (closed and bounded) and
$f\colon K\to\mathbb{R}$ continuous.  Then $f$ attains its maximum and
minimum on~$K$: there exist $a,b\in K$ such that $f(a)\le f(x)\le f(b)$
for all $x\in K$.
\end{theorem}

\begin{proof}
Since $K$ is compact and $f$ is continuous, the image $f(K)$ is a compact
subset of~$\mathbb{R}$ (continuous image of a compact set is compact), hence
closed and bounded.  By completeness of $\mathbb{R}$, $\sup f(K)$ and
$\inf f(K)$ exist and belong to $f(K)$ (since $f(K)$ is closed), giving
points $a,b\in K$ with $f(a)=\inf f(K)$ and $f(b)=\sup f(K)$.
\end{proof}

%% ──────────────────────────────────────────────────────────
\section{Partial Derivatives}
\label{sec:partial-derivatives}
\index{partial derivative}

\begin{definition}[Partial derivative]\label{def:partial-deriv}
Let $f\colon U\subseteq\mathbb{R}^n\to\mathbb{R}$ with $U$ open, and let
$a=(a_1,\dots,a_n)\in U$.  The \emph{partial derivative of $f$ with respect
to $x_i$ at~$a$} is
\[
  \frac{\partial f}{\partial x_i}(a)
  = \lim_{h\to 0}\frac{f(a_1,\dots,a_i+h,\dots,a_n)-f(a)}{h},
\]
provided this limit exists.
\end{definition}

\begin{center}
\begin{tikzpicture}[scale=0.9]
  % Surface cross-section illustration
  \draw[->] (-0.5,0) -- (5,0) node[right]{$x$};
  \draw[->] (0,-0.5) -- (0,4) node[above]{$z$};
  % A curve z = f(x,y_0)
  \draw[thick,blue!70!black,domain=0.3:4.5,samples=80]
    plot (\x, {1.5 + 0.8*sin(60*\x) + 0.15*\x});
  % Tangent line at x=2.5
  \pgfmathsetmacro{\xa}{2.5}
  \pgfmathsetmacro{\ya}{1.5 + 0.8*sin(60*\xa) + 0.15*\xa}
  \pgfmathsetmacro{\sl}{0.8*cos(60*\xa)*3.14159/3 + 0.15}
  \draw[thick, red!70!black, dashed]
    ({\xa-1.2},{\ya - 1.2*\sl}) -- ({\xa+1.2},{\ya + 1.2*\sl});
  \fill ({\xa},{\ya}) circle (1.5pt);
  \node[above right] at ({\xa},{\ya}) {$(a, f(a,y_0))$};
  \node[blue!70!black] at (3.2,0.7) {$z=f(x,y_0)$};
  \node[red!70!black] at (4.5,3.5) {slope $=\dfrac{\partial f}{\partial x}(a,y_0)$};
\end{tikzpicture}
\captionof{figure}{The partial derivative $\partial f/\partial x$ as the slope
of the cross-section $y=y_0$.}
\label{fig:partial-slope}
\end{center}

\begin{theorem}[Existence of partial derivatives does not imply continuity]
\label{thm:partials-not-continuous}
\index{partial derivative!does not imply continuity}
There exists a function $f\colon\mathbb{R}^2\to\mathbb{R}$ such that both
$\dfrac{\partial f}{\partial x}(0,0)$ and
$\dfrac{\partial f}{\partial y}(0,0)$ exist, yet $f$ is not continuous
at~$(0,0)$.
\end{theorem}

\begin{proof}
Define
\[
  f(x,y) = \begin{cases}
    \dfrac{x^2 y}{x^4+y^2} & \text{if }(x,y)\ne(0,0),\\[6pt]
    0 & \text{if }(x,y)=(0,0).
  \end{cases}
\]
\textbf{Partial derivatives at the origin.}\;
$\frac{\partial f}{\partial x}(0,0)
  = \lim_{h\to 0}\frac{f(h,0)-f(0,0)}{h}
  = \lim_{h\to 0}\frac{0}{h}=0$, since $f(h,0)=0$ for all $h\ne 0$.
Similarly, $\frac{\partial f}{\partial y}(0,0)=0$.

\textbf{Discontinuity.}\; Along the curve $y=x^2$,
\[
  f(x,x^2)=\frac{x^2\cdot x^2}{x^4+x^4}=\frac{x^4}{2x^4}=\frac{1}{2}
  \quad\text{for all }x\ne 0.
\]
Therefore $\lim_{x\to 0}f(x,x^2)=\tfrac12\ne 0=f(0,0)$, so $f$ is not
continuous at the origin.
\end{proof}

%% ──────────────────────────────────────────────────────────
\section{Differentiability}
\label{sec:differentiability}
\index{differentiability!in Rn@in $\mathbb{R}^n$}

\begin{tcolorbox}[
  colback=yellow!5!white,
  colframe=orange!70!black,
  fonttitle=\bfseries,
  title={Key Definition: Differentiability},
  breakable
]
\begin{definition}[Differentiability]\label{def:differentiable}
Let $U\subseteq\mathbb{R}^n$ be open, $a\in U$, and
$f\colon U\to\mathbb{R}^m$.  We say $f$ is \emph{differentiable at~$a$}
if there exists a linear map $L\colon\mathbb{R}^n\to\mathbb{R}^m$ such that
\[
  f(a+h) = f(a) + L(h) + \|h\|\,\varepsilon(h),
\]
where $\varepsilon(h)\to 0$ as $h\to 0$.  Equivalently,
\[
  \lim_{h\to 0}\frac{\|f(a+h)-f(a)-L(h)\|}{\|h\|}=0.
\]
The linear map $L$ is unique; it is called the \emph{differential} (or
\emph{total derivative}) of $f$ at~$a$ and is denoted $df_a$ or $Df(a)$.
\index{differential}\index{total derivative}
\end{definition}
\end{tcolorbox}

\begin{theorem}[Differentiable implies continuous]
\label{thm:diff-implies-cont}
\index{differentiability!implies continuity}
If $f$ is differentiable at $a$, then $f$ is continuous at~$a$.
\end{theorem}

\begin{proof}
Write $f(a+h)=f(a)+L(h)+\|h\|\,\varepsilon(h)$ with $\varepsilon(h)\to 0$.
Since $L$ is linear (hence continuous on $\mathbb{R}^n$), $L(h)\to 0$ as
$h\to 0$, and $\|h\|\,\varepsilon(h)\to 0$ as $h\to 0$.  Therefore
$f(a+h)\to f(a)$.
\end{proof}

\begin{theorem}[Differentiable implies partial derivatives exist]
\label{thm:diff-implies-partials}
If $f\colon U\to\mathbb{R}$ is differentiable at $a\in U$, then all partial
derivatives $\frac{\partial f}{\partial x_i}(a)$ exist and
\[
  df_a(h) = \sum_{i=1}^{n}\frac{\partial f}{\partial x_i}(a)\,h_i,
  \qquad h=(h_1,\dots,h_n).
\]
\end{theorem}

\begin{proof}
Let $e_i$ be the $i$-th standard basis vector.  Since $f$ is differentiable,
\[
  f(a+te_i)=f(a)+t\,L(e_i)+|t|\,\varepsilon(te_i)
\]
with $\varepsilon(te_i)\to 0$ as $t\to 0$.  Hence
\[
  \frac{f(a+te_i)-f(a)}{t}=L(e_i)+\frac{|t|}{t}\,\varepsilon(te_i)
  \;\longrightarrow\; L(e_i)
\]
as $t\to 0$.  This shows $\frac{\partial f}{\partial x_i}(a)=L(e_i)$.
Since $L$ is linear and $h=\sum_i h_i e_i$, we get
$L(h)=\sum_i h_i L(e_i)=\sum_i \frac{\partial f}{\partial x_i}(a)\,h_i$.
\end{proof}

\begin{definition}[Jacobian matrix]\label{def:jacobian}
\index{Jacobian matrix}
If $f=(f_1,\dots,f_m)\colon U\subseteq\mathbb{R}^n\to\mathbb{R}^m$ is
differentiable at~$a$, the \emph{Jacobian matrix} of $f$ at $a$ is the
$m\times n$ matrix
\[
  J_f(a) = \begin{pmatrix}
    \dfrac{\partial f_1}{\partial x_1}(a) & \cdots &
    \dfrac{\partial f_1}{\partial x_n}(a) \\[8pt]
    \vdots & \ddots & \vdots \\[4pt]
    \dfrac{\partial f_m}{\partial x_1}(a) & \cdots &
    \dfrac{\partial f_m}{\partial x_n}(a)
  \end{pmatrix}.
\]
Then $df_a(h)=J_f(a)\,h$ for all $h\in\mathbb{R}^n$.
\end{definition}

\begin{theorem}[$C^1$ implies differentiable]
\label{thm:C1-implies-diff}
\index{C1 implies differentiable@$C^1$ implies differentiable}
Let $f\colon U\subseteq\mathbb{R}^n\to\mathbb{R}$ be such that all partial
derivatives $\frac{\partial f}{\partial x_i}$ exist in a neighbourhood of
$a$ and are continuous at~$a$.  Then $f$ is differentiable at~$a$.
\end{theorem}

\begin{proof}[Proof for the case $n=2$]
Write $a=(a_1,a_2)$ and $h=(h_1,h_2)$ with $a+h\in U$.  Then
\begin{align*}
  f(a+h)-f(a)
  &= \bigl[f(a_1+h_1,a_2+h_2)-f(a_1,a_2+h_2)\bigr]
     +\bigl[f(a_1,a_2+h_2)-f(a_1,a_2)\bigr].
\end{align*}
Apply the mean value theorem in one variable to each bracket.  For the first,
the function $t\mapsto f(t,a_2+h_2)$ is differentiable on the segment
$[a_1,a_1+h_1]$, so there exists $\theta_1$ between $a_1$ and $a_1+h_1$
with
\[
  f(a_1+h_1,a_2+h_2)-f(a_1,a_2+h_2)
  = h_1\,\frac{\partial f}{\partial x_1}(\theta_1,a_2+h_2).
\]
Similarly, there exists $\theta_2$ between $a_2$ and $a_2+h_2$ with
\[
  f(a_1,a_2+h_2)-f(a_1,a_2)
  = h_2\,\frac{\partial f}{\partial x_2}(a_1,\theta_2).
\]
Therefore
\begin{align*}
  f(a+h)-f(a)
  &= h_1\,\frac{\partial f}{\partial x_1}(a)
    +h_2\,\frac{\partial f}{\partial x_2}(a) + R(h),
\end{align*}
where
\[
  R(h) = h_1\Bigl[\frac{\partial f}{\partial x_1}(\theta_1,a_2+h_2)
         -\frac{\partial f}{\partial x_1}(a)\Bigr]
        +h_2\Bigl[\frac{\partial f}{\partial x_2}(a_1,\theta_2)
         -\frac{\partial f}{\partial x_2}(a)\Bigr].
\]
Since $(\theta_1,a_2+h_2)\to a$ and $(a_1,\theta_2)\to a$ as $h\to 0$,
the continuity of the partial derivatives at~$a$ gives
\[
  \frac{|R(h)|}{\|h\|}
  \le \Bigl|\frac{\partial f}{\partial x_1}(\theta_1,a_2+h_2)
      -\frac{\partial f}{\partial x_1}(a)\Bigr|
     +\Bigl|\frac{\partial f}{\partial x_2}(a_1,\theta_2)
      -\frac{\partial f}{\partial x_2}(a)\Bigr|
  \;\longrightarrow\; 0,
\]
using $|h_i|\le\|h\|$.  Hence $f$ is differentiable at $a$ with
$df_a(h)=\frac{\partial f}{\partial x_1}(a)h_1
+\frac{\partial f}{\partial x_2}(a)h_2$.
\end{proof}

%% ──────────────────────────────────────────────────────────
\section{The Chain Rule}
\label{sec:chain-rule-multi}
\index{chain rule!multivariable}

\begin{theorem}[Chain rule]
\label{thm:chain-rule}
Let $U\subseteq\mathbb{R}^n$ and $V\subseteq\mathbb{R}^m$ be open.  If
$f\colon U\to V$ is differentiable at~$a$ and $g\colon V\to\mathbb{R}^p$
is differentiable at $f(a)$, then $g\circ f$ is differentiable at~$a$ and
\[
  d(g\circ f)_a = dg_{f(a)}\circ df_a.
\]
In terms of Jacobian matrices, $J_{g\circ f}(a)=J_g(f(a))\cdot J_f(a)$.
\end{theorem}

\begin{proof}
Write $b=f(a)$.  By differentiability,
\begin{align*}
  f(a+h) &= f(a) + df_a(h) + \|h\|\,\varepsilon_1(h), \\
  g(b+k) &= g(b) + dg_b(k) + \|k\|\,\varepsilon_2(k),
\end{align*}
where $\varepsilon_1(h)\to 0$ and $\varepsilon_2(k)\to 0$.  Set
$k=f(a+h)-f(a)=df_a(h)+\|h\|\varepsilon_1(h)$.  Then $\|k\|\le M\|h\|$
for some constant $M>0$ (since $df_a$ is a bounded linear map and
$\|\varepsilon_1(h)\|\to 0$, so $\|k\|\le(\|df_a\|+1)\|h\|$ for $h$ small).

Now
\begin{align*}
  g(f(a+h)) &= g(b) + dg_b(k) + \|k\|\,\varepsilon_2(k) \\
  &= g(b) + dg_b(df_a(h)) + dg_b(\|h\|\varepsilon_1(h))
     + \|k\|\,\varepsilon_2(k).
\end{align*}
The remainder is
\[
  R(h) = \|h\|\,dg_b(\varepsilon_1(h)) + \|k\|\,\varepsilon_2(k).
\]
Since $dg_b$ is continuous and linear, $dg_b(\varepsilon_1(h))\to 0$, so
$\|h\|\,dg_b(\varepsilon_1(h))=o(\|h\|)$.  Also
$\|k\|\,\|\varepsilon_2(k)\|\le M\|h\|\,\|\varepsilon_2(k)\| = o(\|h\|)$
since $\varepsilon_2(k)\to 0$ as $k\to 0$ (and $k\to 0$ as $h\to 0$).
Therefore $\|R(h)\|/\|h\|\to 0$, proving differentiability with
$d(g\circ f)_a = dg_b\circ df_a$.
\end{proof}

%% ──────────────────────────────────────────────────────────
\section{Gradient and Geometric Interpretation}
\label{sec:gradient}
\index{gradient}

\begin{definition}[Gradient]\label{def:gradient}
If $f\colon U\subseteq\mathbb{R}^n\to\mathbb{R}$ is differentiable at~$a$,
the \emph{gradient} of $f$ at $a$ is the vector
\[
  \nabla f(a) = \biggl(\frac{\partial f}{\partial x_1}(a),\dots,
  \frac{\partial f}{\partial x_n}(a)\biggr) \in\mathbb{R}^n.
\]
The differential can then be written $df_a(h)=\langle\nabla f(a),h\rangle$.
\end{definition}

\begin{proposition}[Direction of steepest ascent]
\label{prop:steepest-ascent}
\index{steepest ascent}
Among all unit vectors $v\in\mathbb{R}^n$ with $\|v\|_2=1$, the directional
derivative $D_v f(a)=\langle\nabla f(a),v\rangle$ is maximised when
$v=\nabla f(a)/\|\nabla f(a)\|_2$ (assuming $\nabla f(a)\ne 0$).  The
maximum value is $\|\nabla f(a)\|_2$.
\end{proposition}

\begin{proof}
By Cauchy--Schwarz, $|\langle\nabla f(a),v\rangle|\le\|\nabla f(a)\|_2\|v\|_2
=\|\nabla f(a)\|_2$, with equality when $v$ is parallel to $\nabla f(a)$.
\end{proof}

\begin{remark}
The gradient $\nabla f(a)$ is perpendicular to the level set
$\{x:f(x)=f(a)\}$ at the point~$a$.  This is because for any curve
$\gamma(t)$ lying in the level set with $\gamma(0)=a$, differentiating
$f(\gamma(t))=f(a)$ at $t=0$ gives $\langle\nabla f(a),\gamma'(0)\rangle=0$.
\end{remark}

\begin{center}
\begin{tikzpicture}[scale=0.85]
  % Contour plot with gradient vectors
  \foreach \r in {0.8,1.3,1.8,2.3,2.8} {
    \draw[gray!60, thick] (0,0) circle (\r);
  }
  \node[gray!60] at (2.8,2.2) {$f=c_5$};
  \node[gray!60] at (1.8,1.6) {$f=c_4$};
  % gradient arrows
  \foreach \ang in {0,45,90,135,180,225,270,315} {
    \draw[-{Stealth[length=3mm]}, thick, red!70!black]
      ({1.3*cos(\ang)},{1.3*sin(\ang)}) --
      ({2.0*cos(\ang)},{2.0*sin(\ang)});
  }
  \node[red!70!black] at (3.0,0) {$\nabla f$};
  \node at (0,-3.5) {Gradient vectors point radially outward,};
  \node at (0,-4.0) {perpendicular to level curves.};
\end{tikzpicture}
\captionof{figure}{Gradient vectors on a contour plot of $f(x,y)=x^2+y^2$.}
\label{fig:gradient-contour}
\end{center}

%% ──────────────────────────────────────────────────────────
\section{Higher-Order Partial Derivatives and Schwarz's Theorem}
\label{sec:schwarz}
\index{Schwarz's theorem}
\index{mixed partial derivatives}

\begin{theorem}[Schwarz's theorem --- symmetry of mixed partials]
\label{thm:schwarz}
Let $f\colon U\subseteq\mathbb{R}^2\to\mathbb{R}$ with $U$ open.  Suppose
$\frac{\partial f}{\partial x}$, $\frac{\partial f}{\partial y}$,
$\frac{\partial^2 f}{\partial x\partial y}$, and
$\frac{\partial^2 f}{\partial y\partial x}$ all exist in~$U$ and that
$\frac{\partial^2 f}{\partial x\partial y}$ and
$\frac{\partial^2 f}{\partial y\partial x}$ are continuous at
$(a,b)\in U$.  Then
\[
  \frac{\partial^2 f}{\partial x\partial y}(a,b)
  = \frac{\partial^2 f}{\partial y\partial x}(a,b).
\]
\end{theorem}

\begin{proof}
For $h,k$ small enough, define
\[
  \Phi(h,k)
  = f(a+h,b+k)-f(a+h,b)-f(a,b+k)+f(a,b).
\]
Set $g(t)=f(t,b+k)-f(t,b)$.  Then $\Phi(h,k)=g(a+h)-g(a)$.  By the mean
value theorem, there exists $\theta_1\in(0,1)$ such that
\[
  \Phi(h,k)=h\,g'(a+\theta_1 h)
  = h\Bigl[\frac{\partial f}{\partial x}(a+\theta_1 h,b+k)
    -\frac{\partial f}{\partial x}(a+\theta_1 h,b)\Bigr].
\]
Applying the mean value theorem again (in the $y$-variable), there exists
$\theta_2\in(0,1)$ such that
\[
  \Phi(h,k) = hk\,\frac{\partial^2 f}{\partial y\partial x}
  (a+\theta_1 h,\, b+\theta_2 k).
\]
Similarly, writing $\Phi(h,k)$ using $\tilde{g}(s)=f(a+h,s)-f(a,s)$ and
applying the mean value theorem twice in the other order yields
\[
  \Phi(h,k)=hk\,\frac{\partial^2 f}{\partial x\partial y}
  (a+\tilde\theta_1 h,\, b+\tilde\theta_2 k)
\]
for some $\tilde\theta_1,\tilde\theta_2\in(0,1)$.  Therefore
\[
  \frac{\partial^2 f}{\partial y\partial x}(a+\theta_1 h,b+\theta_2 k)
  = \frac{\partial^2 f}{\partial x\partial y}
    (a+\tilde\theta_1 h, b+\tilde\theta_2 k).
\]
Letting $h,k\to 0$ and using continuity of both mixed partials at $(a,b)$,
we obtain
\[
  \frac{\partial^2 f}{\partial y\partial x}(a,b)
  = \frac{\partial^2 f}{\partial x\partial y}(a,b).  \qedhere
\]
\end{proof}

%% ──────────────────────────────────────────────────────────
\section{Taylor's Formula in Several Variables}
\label{sec:taylor-multi}
\index{Taylor's formula!several variables}

\begin{theorem}[Taylor's formula, order~1]
\label{thm:taylor-1}
Let $f\colon U\subseteq\mathbb{R}^n\to\mathbb{R}$ be $C^1$.  Then for
$a,a+h\in U$,
\[
  f(a+h)=f(a)+\langle\nabla f(a),h\rangle + o(\|h\|).
\]
\end{theorem}

\begin{theorem}[Taylor's formula, order~2]
\label{thm:taylor-2}
Let $f\colon U\subseteq\mathbb{R}^n\to\mathbb{R}$ be $C^2$.  Then
\[
  f(a+h)=f(a)+\sum_{i=1}^n \frac{\partial f}{\partial x_i}(a)\,h_i
  +\frac{1}{2}\sum_{i,j=1}^n
    \frac{\partial^2 f}{\partial x_i\partial x_j}(a)\,h_i h_j
  + o(\|h\|^2).
\]
\end{theorem}

\begin{proof}
Apply the single-variable Taylor formula to $\varphi(t)=f(a+th)$ at $t=0$:
\[
  \varphi(1)=\varphi(0)+\varphi'(0)+\tfrac12\varphi''(0)+o(1).
\]
By the chain rule,
$\varphi'(t)=\sum_i \frac{\partial f}{\partial x_i}(a+th)\,h_i$ and
$\varphi''(t)=\sum_{i,j}\frac{\partial^2 f}{\partial x_i\partial x_j}(a+th)\,h_ih_j$.
Evaluating at $t=0$ yields the stated formula.
\end{proof}

%% ──────────────────────────────────────────────────────────
\section{Extrema of Functions of Several Variables}
\label{sec:extrema-multi}
\index{extrema!several variables}

\begin{definition}[Hessian matrix]\label{def:hessian}
\index{Hessian matrix}
For $f\colon U\subseteq\mathbb{R}^n\to\mathbb{R}$ of class $C^2$, the
\emph{Hessian matrix} at $a\in U$ is the $n\times n$ symmetric matrix
\[
  H_f(a) = \biggl(\frac{\partial^2 f}{\partial x_i\partial x_j}(a)
  \biggr)_{1\le i,j\le n}.
\]
\end{definition}

\begin{theorem}[Necessary condition for an extremum]
\label{thm:necessary-extremum}
\index{critical point}
If $f$ has a local extremum at an interior point $a$ and $f$ is
differentiable at~$a$, then $\nabla f(a)=0$ (i.e.\ $a$ is a
\emph{critical point}).
\end{theorem}

\begin{proof}
For each $i$, the function $t\mapsto f(a+te_i)$ has a local extremum at
$t=0$, so $\frac{\partial f}{\partial x_i}(a)=0$ by the single-variable
result.
\end{proof}

\begin{theorem}[Second-order criterion]
\label{thm:second-order-criterion}
Let $f\colon U\subseteq\mathbb{R}^n\to\mathbb{R}$ be $C^2$ and let
$\nabla f(a)=0$.
\begin{enumerate}[label=(\roman*)]
  \item If $H_f(a)$ is positive definite, then $a$ is a strict local minimum.
  \item If $H_f(a)$ is negative definite, then $a$ is a strict local maximum.
  \item If $H_f(a)$ is indefinite (has both positive and negative eigenvalues),
    then $a$ is a \emph{saddle point}\index{saddle point}.
\end{enumerate}
\end{theorem}

\begin{example}[Classification of critical points in $\mathbb{R}^2$]
\label{ex:critical-points-R2}
For $f\colon\mathbb{R}^2\to\mathbb{R}$ of class $C^2$ with $\nabla f(a)=0$,
set $A=\frac{\partial^2 f}{\partial x^2}(a)$,
$B=\frac{\partial^2 f}{\partial x\partial y}(a)$,
$C=\frac{\partial^2 f}{\partial y^2}(a)$.  Let $\Delta=AC-B^2$.
\begin{itemize}
  \item $\Delta>0,\;A>0$: local minimum.
  \item $\Delta>0,\;A<0$: local maximum.
  \item $\Delta<0$: saddle point.
  \item $\Delta=0$: test inconclusive.
\end{itemize}
\end{example}

\begin{example}
Consider $f(x,y)=x^2-y^2$.  Then $\nabla f=0$ only at the origin.
$A=2$, $B=0$, $C=-2$, $\Delta=2\cdot(-2)-0=-4<0$: saddle point.
\end{example}

\begin{center}
\begin{tikzpicture}[scale=0.7]
  % Saddle point contour sketch
  \begin{scope}
    \draw[->] (-3,0) -- (3,0) node[right]{$x$};
    \draw[->] (0,-3) -- (0,3) node[above]{$y$};
    % Hyperbolas f=c>0
    \draw[blue!70!black, thick, domain=0.71:2.5, samples=50]
      plot (\x, {sqrt(\x*\x - 0.5)}) plot (\x, {-sqrt(\x*\x - 0.5)})
      plot ({-\x}, {sqrt(\x*\x - 0.5)}) plot ({-\x}, {-sqrt(\x*\x - 0.5)});
    % Hyperbolas f=c<0
    \draw[red!70!black, thick, domain=0.71:2.5, samples=50]
      plot ({sqrt(\x*\x - 0.5)}, \x) plot ({-sqrt(\x*\x - 0.5)}, \x)
      plot ({sqrt(\x*\x - 0.5)}, {-\x}) plot ({-sqrt(\x*\x - 0.5)}, {-\x});
    \fill (0,0) circle (2pt);
    \node[below right] at (0.1,-0.2) {saddle};
    \node at (0,-3.8) {$f(x,y)=x^2-y^2$};
  \end{scope}

  \begin{scope}[xshift=8cm]
    \draw[->] (-3,0) -- (3,0) node[right]{$x$};
    \draw[->] (0,-3) -- (0,3) node[above]{$y$};
    % Ellipses f=c>0
    \foreach \r in {0.6,1.1,1.6,2.1,2.6} {
      \draw[blue!70!black, thick] (0,0) circle (\r);
    }
    \fill (0,0) circle (2pt);
    \node[below right] at (0.1,-0.2) {min};
    \node at (0,-3.8) {$f(x,y)=x^2+y^2$};
  \end{scope}
\end{tikzpicture}
\captionof{figure}{Contour plots: saddle point (left) versus local minimum
(right).}
\label{fig:saddle-min}
\end{center}

%% ──────────────────────────────────────────────────────────
\section{Exercises}
\label{sec:ex-ch9}

\begin{exercise}
Show that $\|x\|_2 \le \sqrt{n}\,\|x\|_\infty$ for all $x\in\mathbb{R}^n$.
\end{exercise}

\begin{exercise}
Let $f(x,y)=\frac{x^3y}{x^6+y^2}$ for $(x,y)\ne(0,0)$ and $f(0,0)=0$.
Show that $f$ is not continuous at the origin, even though $f$ is continuous
along every line through the origin.
\end{exercise}

\begin{exercise}
Compute the Jacobian matrix of $f\colon\mathbb{R}^2\to\mathbb{R}^2$ defined
by $f(r,\theta)=(r\cos\theta,r\sin\theta)$.  Show that $\det J_f=r$.
\end{exercise}

\begin{exercise}
Let $f(x,y)=e^{x}\cos y$.  Verify that $\frac{\partial^2 f}{\partial x
\partial y}=\frac{\partial^2 f}{\partial y\partial x}$ everywhere.
\end{exercise}

\begin{exercise}
Let $f(x,y,z)=x^2y+yz^2+zx$.  Compute $\nabla f$ and find all critical
points.
\end{exercise}

\begin{exercise}
Let $f\colon\mathbb{R}^2\to\mathbb{R}$ be defined by
$f(x,y)=x^4+y^4-4xy+1$.  Find all critical points and classify them.
\end{exercise}

\begin{exercise}
Let $g\colon\mathbb{R}^2\to\mathbb{R}$ be differentiable and define
$h(r,\theta)=g(r\cos\theta,r\sin\theta)$.  Express $\frac{\partial h}{\partial r}$
and $\frac{\partial h}{\partial\theta}$ in terms of the partial derivatives of~$g$.
\end{exercise}

\begin{exercise}
Prove that if $f\colon\mathbb{R}^n\to\mathbb{R}$ is differentiable and
$f(tx)=t^k f(x)$ for all $t>0$ (homogeneous of degree~$k$), then
$\langle\nabla f(x),x\rangle=kf(x)$ (Euler's identity).
\end{exercise}

\begin{exercise}
Let $f(x,y)=(x^2+y^2)e^{-(x^2+y^2)}$.  Find and classify all critical
points of~$f$.
\end{exercise}

\begin{exercise}
Write the second-order Taylor expansion of $f(x,y)=\sin(x+2y)$ at
$(0,0)$.
\end{exercise}

\begin{exercise}[$\star\star$]
Show that if $f\colon U\to\mathbb{R}$ is differentiable at $a\in U$ and
attains a local maximum at $a$ subject to the constraint $g(x)=0$
(where $g\colon U\to\mathbb{R}$ is $C^1$ with $\nabla g(a)\ne 0$), then
there exists $\lambda\in\mathbb{R}$ such that
$\nabla f(a)=\lambda\nabla g(a)$.
\end{exercise}

\subsection*{Chapter Summary}

\begin{itemize}
  \item All norms on $\mathbb{R}^n$ are equivalent; topology does not depend
    on the chosen norm.
  \item Existence of partial derivatives alone does not guarantee continuity
    or differentiability.
  \item Differentiability is defined via linear approximation: $f(a+h)
    =f(a)+L(h)+o(\|h\|)$.
  \item $C^1$ implies differentiable; the chain rule composes Jacobians.
  \item The gradient points in the direction of steepest ascent and is
    perpendicular to level sets.
  \item Critical points are classified by the Hessian matrix.
\end{itemize}


%% ══════════════════════════════════════════════════════════
%%  CHAPTER 10
%% ══════════════════════════════════════════════════════════

\chapter{Multiple Integrals and Change of Variables}
\label{ch:multiple-integrals}
\index{multiple integral}

\section*{Introduction}

How does one compute the volume beneath a surface $z=f(x,y)$, the mass of a
solid with variable density, or the centre of mass of a lamina?  These
problems demand integration of functions of several variables.  In this
chapter we define double and triple integrals via Riemann sums, establish
Fubini's theorem (which reduces a multiple integral to iterated
single-variable integrals), and present the change-of-variables formula ---
the multivariable analogue of integration by substitution.

%% ──────────────────────────────────────────────────────────
\section{Double Integrals over Rectangles}
\label{sec:double-int-rect}
\index{double integral!over rectangles}

\begin{definition}{Riemann sum in $\mathbb{R}^2$}{riemann-sum-R2}
Let $R=[a,b]\times[c,d]$ be a closed rectangle and $f\colon R\to\mathbb{R}$
bounded.  A \emph{partition} of $R$ is $P=P_x\times P_y$ where
$P_x\colon a=x_0<x_1<\dots<x_p=b$ and
$P_y\colon c=y_0<y_1<\dots<y_q=d$.  For each sub-rectangle
$R_{ij}=[x_{i-1},x_i]\times[y_{j-1},y_j]$, choose a sample point
$(\xi_{ij},\eta_{ij})\in R_{ij}$.  The \emph{Riemann sum} is
\[
  S(f,P)=\sum_{i=1}^p\sum_{j=1}^q
    f(\xi_{ij},\eta_{ij})\,\Delta x_i\,\Delta y_j,
\]
where $\Delta x_i=x_i-x_{i-1}$ and $\Delta y_j=y_j-y_{j-1}$.
\end{definition}

\begin{definition}[Double integral]\label{def:double-integral}
\index{double integral}
The function $f$ is \emph{Riemann integrable} over $R$ if there exists
$I\in\mathbb{R}$ such that for every $\varepsilon>0$ there exists $\delta>0$
with $|S(f,P)-I|<\varepsilon$ whenever $\|P\|<\delta$ (where $\|P\|$ is the
mesh size), for every choice of sample points.  We write
$I=\iint_R f(x,y)\,dx\,dy$.
\end{definition}

\begin{theorem}[Integrability of continuous functions]
\label{thm:continuous-integrable}
Every continuous function on a closed rectangle is Riemann integrable.
\end{theorem}

%% ──────────────────────────────────────────────────────────
\section{Fubini's Theorem}
\label{sec:fubini}
\index{Fubini's theorem}

\begin{theorem}[Fubini's theorem for continuous functions]
\label{thm:fubini}
Let $R=[a,b]\times[c,d]$ and let $f\colon R\to\mathbb{R}$ be continuous.
Then
\[
  \iint_R f(x,y)\,dx\,dy
  = \int_a^b\!\biggl(\int_c^d f(x,y)\,dy\biggr)dx
  = \int_c^d\!\biggl(\int_a^b f(x,y)\,dx\biggr)dy.
\]
\end{theorem}

\begin{proof}
We prove the first equality; the second follows by symmetry.  Define
$F(x)=\int_c^d f(x,y)\,dy$.  Since $f$ is continuous on the compact set~$R$,
it is uniformly continuous: for every $\varepsilon>0$ there exists $\delta>0$
such that $|f(x,y)-f(x',y')|<\varepsilon$ whenever
$\|(x,y)-(x',y')\|<\delta$.

Let $P=\{x_0,\dots,x_p\}\times\{y_0,\dots,y_q\}$ be a partition of~$R$ with
mesh $\|P\|<\delta$.  For each sub-rectangle $R_{ij}$, the minimum and maximum
of $f$ on $R_{ij}$ satisfy $M_{ij}-m_{ij}<\varepsilon$.

The lower and upper Darboux sums satisfy
\[
  L(f,P) \le \int_a^b F(x)\,dx \le U(f,P).
\]
Indeed, for each $i$,
\[
  \sum_j m_{ij}\,\Delta y_j
  \le \int_c^d f(x,y)\,dy = F(x)
  \le \sum_j M_{ij}\,\Delta y_j
  \qquad\text{for all }x\in[x_{i-1},x_i],
\]
so integrating over $[x_{i-1},x_i]$ and summing over~$i$,
\[
  L(f,P) = \sum_{i,j}m_{ij}\Delta x_i\Delta y_j
  \le \int_a^b F(x)\,dx
  \le \sum_{i,j}M_{ij}\Delta x_i\Delta y_j = U(f,P).
\]
Since $f$ is integrable, $U(f,P)-L(f,P)<\varepsilon\,|R|$ (where $|R|$
is the area of~$R$) for $\|P\|$ small enough.  Because the double integral
$\iint_R f$ also lies between $L(f,P)$ and $U(f,P)$, we conclude
\[
  \biggl|\iint_R f\,dx\,dy - \int_a^b F(x)\,dx\biggr|
  \le U(f,P)-L(f,P)<\varepsilon\,|R|.
\]
Since $\varepsilon$ is arbitrary, the two quantities are equal.
\end{proof}

%% ──────────────────────────────────────────────────────────
\section{Integration over Bounded Domains}
\label{sec:bounded-domains}
\index{double integral!over bounded domains}

\begin{definition}[Type I and Type II domains]\label{def:type-I-II}
\index{Type I domain}\index{Type II domain}
A \emph{Type~I} domain has the form
$D=\{(x,y):a\le x\le b,\;\varphi_1(x)\le y\le\varphi_2(x)\}$
for continuous $\varphi_1,\varphi_2$.
A \emph{Type~II} domain has the form
$D=\{(x,y):c\le y\le d,\;\psi_1(y)\le x\le\psi_2(y)\}$.
\end{definition}

\begin{center}
\begin{tikzpicture}[scale=1.1]
  % Type I
  \begin{scope}
    \draw[->] (-0.3,0) -- (4,0) node[right]{$x$};
    \draw[->] (0,-0.3) -- (0,3.2) node[above]{$y$};
    \draw[thick, blue!70!black, domain=0.5:3.3, samples=50]
      plot (\x, {0.3+0.2*sin(120*\x)});
    \draw[thick, red!70!black, domain=0.5:3.3, samples=50]
      plot (\x, {2.5-0.3*cos(80*\x)});
    \draw[dashed] (0.5,0) -- (0.5,{0.3+0.2*sin(60)});
    \draw[dashed] (3.3,0) -- (3.3,{0.3+0.2*sin(396)});
    \node at (0.5,-0.25) {$a$};
    \node at (3.3,-0.25) {$b$};
    \node[blue!70!black] at (4.0,0.4) {$y=\varphi_1(x)$};
    \node[red!70!black] at (4.0,2.3) {$y=\varphi_2(x)$};
    \node at (1.9,1.4) {$D$};
    \node at (1.9,-0.8) {Type I domain};
    % shading
    \fill[blue!10, domain=0.5:3.3, samples=50]
      plot (\x, {0.3+0.2*sin(120*\x)}) --
      plot[domain=3.3:0.5] (\x, {2.5-0.3*cos(80*\x)}) -- cycle;
  \end{scope}

  \begin{scope}[xshift=7cm]
    \draw[->] (-0.3,0) -- (4,0) node[right]{$x$};
    \draw[->] (0,-0.3) -- (0,3.2) node[above]{$y$};
    \draw[thick, blue!70!black, domain=0.5:2.8, samples=50]
      plot ({0.4+0.3*cos(100*\x)}, \x);
    \draw[thick, red!70!black, domain=0.5:2.8, samples=50]
      plot ({3.0-0.2*sin(130*\x)}, \x);
    \draw[dashed] (0,0.5) -- ({0.4+0.3*cos(50)},0.5);
    \draw[dashed] (0,2.8) -- ({0.4+0.3*cos(280)},2.8);
    \node at (-0.3,0.5) {$c$};
    \node at (-0.3,2.8) {$d$};
    \node at (1.9,1.6) {$D$};
    \node at (1.9,-0.8) {Type II domain};
    \fill[red!10, domain=0.5:2.8, samples=50]
      plot ({0.4+0.3*cos(100*\x)}, \x) --
      plot[domain=2.8:0.5] ({3.0-0.2*sin(130*\x)}, \x) -- cycle;
  \end{scope}
\end{tikzpicture}
\captionof{figure}{Type~I (left) and Type~II (right) integration domains.}
\label{fig:type-I-II}
\end{center}

For a Type~I domain,
\[
  \iint_D f\,dA = \int_a^b\int_{\varphi_1(x)}^{\varphi_2(x)}
  f(x,y)\,dy\,dx.
\]

\begin{example}[Integration over a triangle]
\label{ex:triangle-integral}
Compute $\iint_D (x+y)\,dA$ where $D$ is the triangle with vertices
$(0,0)$, $(1,0)$, $(0,1)$.

The domain is of Type~I: $0\le x\le 1$, $0\le y\le 1-x$.  Thus
\begin{align*}
  \iint_D (x+y)\,dA
  &= \int_0^1\int_0^{1-x}(x+y)\,dy\,dx
  = \int_0^1\Bigl[xy+\frac{y^2}{2}\Bigr]_0^{1-x}dx \\
  &= \int_0^1\Bigl(x(1-x)+\frac{(1-x)^2}{2}\Bigr)dx
  = \int_0^1\frac{1-x^2}{2}\,dx \\
  &= \frac{1}{2}\Bigl[x-\frac{x^3}{3}\Bigr]_0^1
  = \frac{1}{2}\cdot\frac{2}{3}=\frac{1}{3}.
\end{align*}
\end{example}

%% ──────────────────────────────────────────────────────────
\section{Change of Variables}
\label{sec:change-of-var}
\index{change of variables!multiple integrals}

\begin{theorem}[Change of variables formula]
\label{thm:change-of-var}
\index{Jacobian determinant}
Let $\Phi\colon U\to V$ be a $C^1$ diffeomorphism between open sets
$U,V\subseteq\mathbb{R}^n$, and let $f\colon V\to\mathbb{R}$ be continuous
with compact support in~$V$.  Then
\[
  \int_V f(y)\,dy = \int_U f(\Phi(x))\,|\det J_\Phi(x)|\,dx.
\]
\end{theorem}

\begin{remark}
The factor $|\det J_\Phi(x)|$ accounts for how the mapping $\Phi$ locally
stretches or compresses volumes.
\end{remark}

%% ──────────────────────────────────────────────────────────
\section{Polar Coordinates}
\label{sec:polar-coords}
\index{polar coordinates}

The map $\Phi(r,\theta)=(r\cos\theta,r\sin\theta)$ has Jacobian matrix
\[
  J_\Phi = \begin{pmatrix}
    \cos\theta & -r\sin\theta \\
    \sin\theta & r\cos\theta
  \end{pmatrix},
  \qquad \det J_\Phi = r.
\]
Therefore $dA = r\,dr\,d\theta$.

\begin{center}
\begin{tikzpicture}[scale=0.95]
  % Polar grid
  \draw[->] (-3.2,0) -- (3.5,0) node[right]{$x$};
  \draw[->] (0,-3.2) -- (0,3.5) node[above]{$y$};
  \foreach \r in {0.5,1.0,...,3.0} {
    \draw[gray!50] (0,0) circle (\r);
  }
  \foreach \a in {0,30,...,330} {
    \draw[gray!50] (0,0) -- ({3*cos(\a)},{3*sin(\a)});
  }
  \draw[->, thick, red!70!black] (0,0) -- ({2*cos(40)},{2*sin(40)})
    node[midway, above left]{$r$};
  \draw[thick, red!70!black] (0.7,0) arc (0:40:0.7);
  \node[red!70!black] at (1.0,0.3) {$\theta$};
\end{tikzpicture}
\captionof{figure}{Polar coordinate grid: concentric circles $r=\text{const}$
and rays $\theta=\text{const}$.}
\label{fig:polar-grid}
\end{center}

\begin{example}[Area of a disk]
\label{ex:area-disk}
\[
  \text{Area}(D_R) = \iint_{D_R}dA
  = \int_0^{2\pi}\int_0^R r\,dr\,d\theta
  = 2\pi\cdot\frac{R^2}{2}=\pi R^2.
\]
\end{example}

\begin{theorem}[The Gaussian integral]
\label{thm:gaussian}
\index{Gaussian integral}
\[
  \int_{-\infty}^{+\infty}e^{-x^2}\,dx = \sqrt{\pi}.
\]
\end{theorem}

\begin{proof}
Let $I=\int_{-\infty}^{+\infty}e^{-x^2}\,dx$.  Then
\[
  I^2 = \int_{-\infty}^{+\infty}e^{-x^2}dx
        \int_{-\infty}^{+\infty}e^{-y^2}dy
  = \iint_{\mathbb{R}^2}e^{-(x^2+y^2)}\,dx\,dy.
\]
Switching to polar coordinates $(x,y)=(r\cos\theta,r\sin\theta)$ with
$dA=r\,dr\,d\theta$:
\[
  I^2 = \int_0^{2\pi}\int_0^{\infty}e^{-r^2}\,r\,dr\,d\theta
  = 2\pi\int_0^{\infty}r\,e^{-r^2}\,dr.
\]
Substituting $u=r^2$, $du=2r\,dr$:
\[
  \int_0^{\infty}r\,e^{-r^2}\,dr = \frac{1}{2}\int_0^{\infty}e^{-u}\,du
  = \frac{1}{2}.
\]
Therefore $I^2=2\pi\cdot\frac{1}{2}=\pi$, and since $I>0$, we conclude
$I=\sqrt{\pi}$.
\end{proof}

%% ──────────────────────────────────────────────────────────
\section{Triple Integrals}
\label{sec:triple-integrals}
\index{triple integral}

Fubini's theorem extends naturally:

\begin{theorem}[Fubini in $\mathbb{R}^3$]
\label{thm:fubini-3d}
If $f\colon[a_1,b_1]\times[a_2,b_2]\times[a_3,b_3]\to\mathbb{R}$ is
continuous, then
\[
  \iiint f\,dx\,dy\,dz
  = \int_{a_1}^{b_1}\int_{a_2}^{b_2}\int_{a_3}^{b_3}
    f(x,y,z)\,dz\,dy\,dx,
\]
and the order of integration may be permuted freely.
\end{theorem}

%% ──────────────────────────────────────────────────────────
\section{Cylindrical and Spherical Coordinates}
\label{sec:cyl-sph}
\index{cylindrical coordinates}\index{spherical coordinates}

\subsection{Cylindrical coordinates}

The map $(r,\theta,z)\mapsto(r\cos\theta,r\sin\theta,z)$ has Jacobian
determinant~$r$.  Hence $dV=r\,dr\,d\theta\,dz$.

\subsection{Spherical coordinates}

\begin{definition}[Spherical coordinates]\label{def:spherical-coords}
The map $\Phi(\rho,\theta,\varphi)
=(\rho\sin\varphi\cos\theta,\;\rho\sin\varphi\sin\theta,\;\rho\cos\varphi)$,
where $\rho\ge 0$, $\theta\in[0,2\pi)$, $\varphi\in[0,\pi]$.
\end{definition}

The Jacobian matrix is
\[
  J_\Phi = \begin{pmatrix}
    \sin\varphi\cos\theta & -\rho\sin\varphi\sin\theta
      & \rho\cos\varphi\cos\theta \\
    \sin\varphi\sin\theta & \rho\sin\varphi\cos\theta
      & \rho\cos\varphi\sin\theta \\
    \cos\varphi & 0 & -\rho\sin\varphi
  \end{pmatrix}.
\]
A direct computation gives $\det J_\Phi = -\rho^2\sin\varphi$, so
\[
  |\det J_\Phi| = \rho^2\sin\varphi,
  \qquad dV = \rho^2\sin\varphi\,d\rho\,d\theta\,d\varphi.
\]

\begin{center}
\tdplotsetmaincoords{70}{120}
\begin{tikzpicture}[scale=1.6, tdplot_main_coords]
  % Spherical coordinates illustration
  \draw[->] (0,0,0) -- (2.5,0,0) node[right]{$x$};
  \draw[->] (0,0,0) -- (0,2.5,0) node[above right]{$y$};
  \draw[->] (0,0,0) -- (0,0,2.5) node[above]{$z$};
  % point
  \pgfmathsetmacro{\rr}{2}
  \pgfmathsetmacro{\tth}{55}
  \pgfmathsetmacro{\pph}{40}
  \coordinate (P) at
    ({\rr*sin(\pph)*cos(\tth)},{\rr*sin(\pph)*sin(\tth)},{\rr*cos(\pph)});
  \draw[thick, red!70!black, ->] (0,0,0) -- (P) node[above right]{$P$};
  % projection
  \coordinate (Pxy) at
    ({\rr*sin(\pph)*cos(\tth)},{\rr*sin(\pph)*sin(\tth)},0);
  \draw[dashed] (P) -- (Pxy);
  \draw[dashed] (0,0,0) -- (Pxy);
  % angles
  \draw[blue!70!black] (0.5,0,0) arc (0:\tth:0.5);
  \node[blue!70!black] at (0.55,0.25,0) {$\theta$};
  \draw[green!50!black] (0,0,0.6) -- ({0.6*sin(20)*cos(\tth)},{0.6*sin(20)*sin(\tth)},{0.6*cos(20)});
  \node[green!50!black] at (0.15,0.05,0.75) {$\varphi$};
  \node[red!70!black] at (0.7,0.5,1.0) {$\rho$};
\end{tikzpicture}
\captionof{figure}{Spherical coordinates $(\rho,\theta,\varphi)$.}
\label{fig:spherical}
\end{center}

\begin{example}[Volume of the ball $B_R$]
\label{ex:volume-ball}
\begin{align*}
  \text{Vol}(B_R)
  &= \int_0^{2\pi}\!\int_0^{\pi}\!\int_0^{R}
     \rho^2\sin\varphi\,d\rho\,d\varphi\,d\theta \\
  &= \Bigl(\int_0^{2\pi}d\theta\Bigr)
     \Bigl(\int_0^{\pi}\sin\varphi\,d\varphi\Bigr)
     \Bigl(\int_0^{R}\rho^2\,d\rho\Bigr) \\
  &= 2\pi\cdot 2\cdot\frac{R^3}{3} = \frac{4\pi R^3}{3}.
\end{align*}
\end{example}

\begin{example}[Moment of inertia of a solid ball]
\label{ex:moment-inertia-ball}
\index{moment of inertia}
The moment of inertia of a homogeneous ball of mass $M$ and radius $R$
about a diameter (say the $z$-axis) is
\[
  I_z = \frac{\rho_0}{1}\iiint_{B_R}(x^2+y^2)\,dV,
\]
where $\rho_0=3M/(4\pi R^3)$ is the density.  In spherical coordinates,
$x^2+y^2=\rho^2\sin^2\varphi$, so
\begin{align*}
  I_z &= \rho_0\int_0^{2\pi}\!\int_0^{\pi}\!\int_0^R
    \rho^2\sin^2\varphi\cdot\rho^2\sin\varphi\,d\rho\,d\varphi\,d\theta \\
  &= \rho_0\cdot 2\pi\cdot\int_0^{\pi}\sin^3\varphi\,d\varphi
    \cdot\int_0^R\rho^4\,d\rho \\
  &= \rho_0\cdot 2\pi\cdot\frac{4}{3}\cdot\frac{R^5}{5}
  = \frac{3M}{4\pi R^3}\cdot\frac{8\pi R^5}{15}
  = \frac{2}{5}MR^2.
\end{align*}
\end{example}

%% ──────────────────────────────────────────────────────────
\section{Exercises}
\label{sec:ex-ch10}

\begin{exercise}
Compute $\iint_R xy\,dA$ where $R=[0,1]\times[0,2]$.
\end{exercise}

\begin{exercise}
Compute $\iint_D e^{x+y}\,dA$ where
$D=\{(x,y):0\le x\le 1,\;0\le y\le x\}$.
\end{exercise}

\begin{exercise}
Switch the order of integration:
$\int_0^1\int_0^{\sqrt{x}}f(x,y)\,dy\,dx$.
\end{exercise}

\begin{exercise}
Use polar coordinates to compute $\iint_D\sqrt{x^2+y^2}\,dA$ over the disk
$D=\{x^2+y^2\le 4\}$.
\end{exercise}

\begin{exercise}
Compute $\iint_D \frac{1}{(1+x^2+y^2)^2}\,dA$ over $\mathbb{R}^2$.
\end{exercise}

\begin{exercise}
Compute $\iiint_E z\,dV$ where $E$ is the tetrahedron with vertices
$(0,0,0)$, $(1,0,0)$, $(0,1,0)$, $(0,0,1)$.
\end{exercise}

\begin{exercise}
Use cylindrical coordinates to compute $\iiint_E dV$ where
$E=\{x^2+y^2\le 1,\;0\le z\le 2-x^2-y^2\}$.
\end{exercise}

\begin{exercise}
Compute the volume enclosed by the sphere $\rho=2$ and the cone
$\varphi=\pi/3$ (using spherical coordinates).
\end{exercise}

\begin{exercise}
Find the centre of mass of the region bounded by $y=x^2$ and $y=x$
with uniform density.
\end{exercise}

\begin{exercise}
Show that $\int_0^{\infty}e^{-x^2}\cos(2bx)\,dx=\frac{\sqrt{\pi}}{2}e^{-b^2}$
for all $b\in\mathbb{R}$.
\textit{Hint:} differentiate $I(b)=\int_0^{\infty}e^{-(x^2+2ibx)}\,dx$
under the integral sign.
\end{exercise}

\begin{exercise}[$\star\star$]
Compute $\displaystyle\int_0^1\!\int_0^1\frac{1}{1-xy}\,dx\,dy$ and deduce that
$\displaystyle\sum_{n=1}^{\infty}\frac{1}{n^2}=\frac{\pi^2}{6}$.
\end{exercise}

\subsection*{Chapter Summary}

\begin{itemize}
  \item Double and triple integrals are defined via Riemann sums over
    partitions of rectangles/boxes.
  \item Fubini's theorem reduces multiple integrals to iterated
    single-variable integrals.
  \item The change-of-variables formula introduces the Jacobian determinant
    as a ``volume scaling factor.''
  \item Polar, cylindrical, and spherical coordinates simplify integration
    over regions with circular or spherical symmetry.
  \item The Gaussian integral $\int e^{-x^2}dx=\sqrt\pi$ is elegantly
    proved via polar coordinates.
\end{itemize}


%% ══════════════════════════════════════════════════════════
%%  CHAPTER 11
%% ══════════════════════════════════════════════════════════

\chapter{Curves and Line Integrals}
\label{ch:curves-line-integrals}
\index{line integral}

\section*{Introduction}

A particle moves through space along a trajectory; a force field does work
on the particle as it travels.  Computing this work requires integrating a
vector field along a curve --- a \emph{line integral}.  This chapter
introduces parametric curves, arc length, and both scalar and vector line
integrals.  We then study conservative (gradient) fields, for which the line
integral depends only on the endpoints, and conclude with Green's theorem,
which relates a line integral around a closed curve to a double integral over
the enclosed region.

%% ──────────────────────────────────────────────────────────
\section{Parametric Curves}
\label{sec:parametric-curves}
\index{parametric curve}

\begin{definition}[Parametric curve]\label{def:parametric-curve}
A \emph{parametric curve} in $\mathbb{R}^n$ is a continuous map
$\gamma\colon[a,b]\to\mathbb{R}^n$.  The image $\gamma([a,b])$ is called
the \emph{trace} (or \emph{support}) of the curve.
\end{definition}

\begin{definition}[Regular curve, piecewise $C^1$ curve]\label{def:regular-curve}
\index{regular curve}\index{piecewise C1 curve@piecewise $C^1$ curve}
A parametric curve $\gamma\colon[a,b]\to\mathbb{R}^n$ of class $C^1$ is
\emph{regular} if $\gamma'(t)\ne 0$ for all $t\in[a,b]$.  It is
\emph{piecewise $C^1$} if there exists a partition
$a=t_0<t_1<\dots<t_N=b$ such that $\gamma|_{[t_{k-1},t_k]}$ is $C^1$ for
each~$k$.
\end{definition}

\begin{definition}[Tangent vector and tangent line]\label{def:tangent-vector}
\index{tangent vector}\index{tangent line}
If $\gamma$ is differentiable at $t_0$, the \emph{tangent vector} is
$\gamma'(t_0)$.  The \emph{tangent line} at $\gamma(t_0)$ is
\[
  \ell(s) = \gamma(t_0) + s\,\gamma'(t_0), \qquad s\in\mathbb{R}.
\]
\end{definition}

\begin{center}
\begin{tikzpicture}[scale=1.2, decoration={markings,
  mark=between positions 0.15 and 0.85 step 0.35
  with {\arrow{Stealth[length=3mm]}}}]
  % A curve
  \draw[thick, blue!70!black, postaction={decorate},
    domain=0.3:3.5, samples=80]
    plot (\x, {1.0 + 0.8*sin(100*\x) + 0.1*\x});
  % Tangent vectors at selected points
  \foreach \t/\lab in {1.0/t_1, 2.0/t_2, 3.0/t_3} {
    \pgfmathsetmacro{\yy}{1.0 + 0.8*sin(100*\t) + 0.1*\t}
    \pgfmathsetmacro{\dy}{0.8*cos(100*\t)*100*3.14159/180 + 0.1}
    \pgfmathsetmacro{\len}{0.6/sqrt(1+\dy*\dy)}
    \fill[red!70!black] (\t,\yy) circle (1.5pt);
    \draw[-{Stealth[length=2.5mm]}, thick, red!70!black]
      (\t,\yy) -- ({\t+\len},{\yy+\len*\dy});
    \node[below, font=\small] at (\t,{\yy-0.2}) {$\gamma(\lab)$};
  }
  \node[red!70!black] at (4.0,2.0) {$\gamma'(t)$};
\end{tikzpicture}
\captionof{figure}{A parametric curve with tangent vectors at several points.}
\label{fig:tangent-vectors}
\end{center}

%% ──────────────────────────────────────────────────────────
\section{Examples of Curves}
\label{sec:curve-examples}

\begin{example}[Circle]
\label{ex:circle}
$\gamma(t)=(\cos t,\sin t)$, $t\in[0,2\pi]$.  This is a regular $C^\infty$
curve; $\gamma'(t)=(-\sin t,\cos t)$, $\|\gamma'(t)\|=1$.
\end{example}

\begin{example}[Ellipse]
\label{ex:ellipse}
$\gamma(t)=(a\cos t,b\sin t)$, $t\in[0,2\pi]$, with $a,b>0$.
$\gamma'(t)=(-a\sin t,b\cos t)\ne 0$ since $a,b>0$.
\end{example}

\begin{example}[Cycloid]
\label{ex:cycloid}
\index{cycloid}
$\gamma(t)=(t-\sin t,1-\cos t)$, $t\in[0,2\pi]$.  Here
$\gamma'(t)=(1-\cos t,\sin t)$, which vanishes at $t=0$ and $t=2\pi$
(cusps).
\end{example}

\begin{example}[Helix]
\label{ex:helix}
\index{helix}
$\gamma(t)=(\cos t,\sin t,t)$, $t\in[0,4\pi]$.  This is a regular curve in
$\mathbb{R}^3$; $\gamma'(t)=(-\sin t,\cos t,1)$,
$\|\gamma'(t)\|=\sqrt{2}$.
\end{example}

%% ──────────────────────────────────────────────────────────
\section{Arc Length}
\label{sec:arc-length}
\index{arc length}

\begin{definition}[Arc length]\label{def:arc-length}
Let $\gamma\colon[a,b]\to\mathbb{R}^n$ be a continuous curve.  For a
partition $P\colon a=t_0<t_1<\dots<t_N=b$, define the \emph{polygonal
length}
\[
  \ell(\gamma,P)=\sum_{k=1}^N\|\gamma(t_k)-\gamma(t_{k-1})\|.
\]
The \emph{arc length} (or \emph{length}) of $\gamma$ is
\[
  L(\gamma)=\sup_P \ell(\gamma,P).
\]
If $L(\gamma)<\infty$, the curve is called \emph{rectifiable}.
\end{definition}

\begin{theorem}[Arc length formula]
\label{thm:arc-length-formula}
If $\gamma\colon[a,b]\to\mathbb{R}^n$ is $C^1$, then $\gamma$ is rectifiable
and
\[
  L(\gamma)=\int_a^b\|\gamma'(t)\|\,dt.
\]
\end{theorem}

\begin{proof}
\textbf{Upper bound.}\; For any partition $P=\{t_0,\dots,t_N\}$,
\[
  \|\gamma(t_k)-\gamma(t_{k-1})\|
  = \Bigl\|\int_{t_{k-1}}^{t_k}\gamma'(t)\,dt\Bigr\|
  \le \int_{t_{k-1}}^{t_k}\|\gamma'(t)\|\,dt.
\]
Summing over $k$: $\ell(\gamma,P)\le\int_a^b\|\gamma'(t)\|\,dt$.
Taking the supremum over~$P$,
$L(\gamma)\le\int_a^b\|\gamma'(t)\|\,dt$.

\textbf{Lower bound.}\; Since $\gamma'$ is continuous on $[a,b]$, it is
uniformly continuous.  Let $\varepsilon>0$; choose $\delta>0$ such that
$|t-s|<\delta$ implies $\|\gamma'(t)-\gamma'(s)\|<\varepsilon$.  Take a
partition with $\|P\|<\delta$.  For $t\in[t_{k-1},t_k]$,
\[
  \gamma'(t) = \gamma'(t_{k-1}) + (\gamma'(t)-\gamma'(t_{k-1})),
\]
so
\[
  \int_{t_{k-1}}^{t_k}\gamma'(t)\,dt
  = \gamma'(t_{k-1})\,(t_k-t_{k-1}) + R_k,
  \quad \|R_k\|\le\varepsilon(t_k-t_{k-1}).
\]
Therefore
\begin{align*}
  \|\gamma(t_k)-\gamma(t_{k-1})\|
  &\ge \|\gamma'(t_{k-1})\|\,(t_k-t_{k-1}) - \|R_k\| \\
  &\ge \|\gamma'(t_{k-1})\|\,(t_k-t_{k-1}) - \varepsilon(t_k-t_{k-1}).
\end{align*}
Summing,
\[
  \ell(\gamma,P)
  \ge \sum_k \|\gamma'(t_{k-1})\|\,(t_k-t_{k-1}) - \varepsilon(b-a).
\]
The sum on the right is a Riemann sum for $\int_a^b\|\gamma'(t)\|\,dt$, so
for $\|P\|$ small enough it differs from the integral by less than
$\varepsilon$.  Thus
$L(\gamma)\ge\int_a^b\|\gamma'(t)\|\,dt-2\varepsilon(b-a+1)$.
Since $\varepsilon$ is arbitrary, $L(\gamma)\ge\int_a^b\|\gamma'(t)\|\,dt$.
\end{proof}

\begin{example}[Length of the helix]
\label{ex:helix-length}
For $\gamma(t)=(\cos t,\sin t,t)$ on $[0,2\pi]$,
$\|\gamma'(t)\|=\sqrt{\sin^2 t+\cos^2 t+1}=\sqrt{2}$, so
$L=\int_0^{2\pi}\sqrt{2}\,dt=2\pi\sqrt{2}$.
\end{example}

%% ──────────────────────────────────────────────────────────
\section{Arc Length Parameterisation}
\label{sec:arc-length-param}
\index{arc length parameterisation}

\begin{definition}[Arc length parameter]\label{def:arc-length-param}
Given a regular $C^1$ curve $\gamma\colon[a,b]\to\mathbb{R}^n$, the
\emph{arc length function} is
$s(t)=\int_a^t\|\gamma'(\tau)\|\,d\tau$.
Since $s'(t)=\|\gamma'(t)\|>0$, the function $s$ is strictly increasing with
a $C^1$ inverse $t=t(s)$.  The reparameterisation
$\tilde\gamma(s)=\gamma(t(s))$ satisfies $\|\tilde\gamma'(s)\|=1$
for all~$s$.
\end{definition}

%% ──────────────────────────────────────────────────────────
\section{Scalar Line Integrals}
\label{sec:scalar-line-integral}
\index{line integral!scalar}

\begin{definition}[Scalar line integral]\label{def:scalar-line-integral}
Let $\gamma\colon[a,b]\to\mathbb{R}^n$ be piecewise $C^1$ and
$f\colon\gamma([a,b])\to\mathbb{R}$ continuous.  The \emph{scalar line
integral} of $f$ along $\gamma$ is
\[
  \int_\gamma f\,ds = \int_a^b f(\gamma(t))\,\|\gamma'(t)\|\,dt.
\]
\end{definition}

\begin{example}[Mass of a wire]
\label{ex:mass-wire}
\index{mass of a wire}
A wire has the shape of the helix $\gamma(t)=(\cos t,\sin t,t)$,
$t\in[0,2\pi]$, with linear density $\rho(x,y,z)=z$.  Its mass is
\[
  M = \int_\gamma \rho\,ds
  = \int_0^{2\pi}t\cdot\sqrt{2}\,dt
  = \sqrt{2}\cdot\frac{(2\pi)^2}{2}
  = 2\pi^2\sqrt{2}.
\]
\end{example}

%% ──────────────────────────────────────────────────────────
\section{Vector Line Integrals (Work)}
\label{sec:vector-line-integral}
\index{line integral!vector}\index{work}

\begin{definition}[Vector line integral]\label{def:vector-line-integral}
Let $F\colon U\subseteq\mathbb{R}^n\to\mathbb{R}^n$ be a continuous vector
field and $\gamma\colon[a,b]\to U$ a piecewise $C^1$ curve.  The
\emph{line integral of $F$ along $\gamma$} is
\[
  \int_\gamma F\cdot d\gamma
  = \int_a^b \langle F(\gamma(t)),\gamma'(t)\rangle\,dt.
\]
\end{definition}

\begin{remark}
Physically, if $F$ is a force field and $\gamma$ is the trajectory of a
particle, then $\int_\gamma F\cdot d\gamma$ is the \emph{work} done by~$F$
on the particle.
\end{remark}

\begin{center}
\begin{tikzpicture}[scale=1.1, decoration={markings,
  mark=between positions 0.1 and 0.9 step 0.2
  with {\arrow{Stealth[length=3mm]}}}]
  % A curve
  \draw[thick, blue!70!black, postaction={decorate},
    domain=0:3.14, samples=60]
    plot ({2*cos(\x r)+2.5},{1.5*sin(\x r)+1.5});
  % Force vectors at several points along curve
  \foreach \t in {0.3, 0.9, 1.5, 2.1, 2.7} {
    \pgfmathsetmacro{\px}{2*cos(\t r)+2.5}
    \pgfmathsetmacro{\py}{1.5*sin(\t r)+1.5}
    \draw[-{Stealth[length=2mm]}, thick, red!70!black]
      (\px,\py) -- ({\px+0.4},{\py+0.25});
  }
  \node[blue!70!black] at (4.8,1.5) {$\gamma$};
  \node[red!70!black] at (5.2,2.5) {$F$};
  \node at (2.5,-0.3) {Work $=\displaystyle\int_\gamma F\cdot d\gamma$};
\end{tikzpicture}
\captionof{figure}{A force field $F$ evaluated along a curve $\gamma$.
The line integral gives the total work.}
\label{fig:force-field-curve}
\end{center}

%% ──────────────────────────────────────────────────────────
\section{Gradient Fields and Potentials}
\label{sec:gradient-fields}
\index{gradient field}\index{potential}

\begin{definition}[Conservative field]\label{def:conservative-field}
\index{conservative field}
A vector field $F\colon U\to\mathbb{R}^n$ is \emph{conservative} (or a
\emph{gradient field}) if there exists a $C^1$ function
$V\colon U\to\mathbb{R}$ such that $F=\nabla V$.  The function $V$ is called
a \emph{potential} for~$F$.
\end{definition}

\begin{theorem}[Fundamental theorem for line integrals]
\label{thm:ftli}
\index{fundamental theorem for line integrals}
If $F=\nabla V$ on an open set $U$ and $\gamma\colon[a,b]\to U$ is
piecewise $C^1$, then
\[
  \int_\gamma F\cdot d\gamma = V(\gamma(b))-V(\gamma(a)).
\]
In particular, $\int_\gamma F\cdot d\gamma$ depends only on the endpoints,
and $\oint_\gamma F\cdot d\gamma=0$ for every closed curve.
\end{theorem}

\begin{proof}
By the chain rule,
\[
  \frac{d}{dt}\bigl[V(\gamma(t))\bigr]
  = \langle\nabla V(\gamma(t)),\gamma'(t)\rangle
  = \langle F(\gamma(t)),\gamma'(t)\rangle.
\]
Integrating from $a$ to $b$:
\[
  \int_a^b \langle F(\gamma(t)),\gamma'(t)\rangle\,dt
  = V(\gamma(b))-V(\gamma(a)).  \qedhere
\]
\end{proof}

\begin{proposition}[Necessary condition for a gradient field in $\mathbb{R}^2$]
\label{prop:necessary-gradient}
Let $F=(P,Q)\colon U\subseteq\mathbb{R}^2\to\mathbb{R}^2$ be $C^1$ with $U$
open.  If $F=\nabla V$ for some $C^2$ function $V$, then
\[
  \frac{\partial P}{\partial y}=\frac{\partial Q}{\partial x}
  \qquad\text{on } U.
\]
\end{proposition}

\begin{proof}
$P=\frac{\partial V}{\partial x}$ and $Q=\frac{\partial V}{\partial y}$, so
$\frac{\partial P}{\partial y}=\frac{\partial^2 V}{\partial y\partial x}
=\frac{\partial^2 V}{\partial x\partial y}=\frac{\partial Q}{\partial x}$
by Schwarz's theorem (Theorem~\ref{thm:schwarz}).
\end{proof}

\begin{remark}
When $U$ is \emph{simply connected}, the condition
$\frac{\partial P}{\partial y}=\frac{\partial Q}{\partial x}$ is also
sufficient for $F$ to be a gradient field.  On domains that are not simply
connected, this can fail; see the counterexample in
Section~\ref{sec:closed-not-exact}.
\end{remark}

%% ──────────────────────────────────────────────────────────
\section{Green's Theorem}
\label{sec:greens-theorem}
\index{Green's theorem}

\begin{theorem}[Green's theorem]
\label{thm:green}
Let $D\subseteq\mathbb{R}^2$ be a bounded domain whose boundary $\partial D$
is a piecewise $C^1$ simple closed curve, oriented counterclockwise.  Let
$P,Q\colon\overline{D}\to\mathbb{R}$ be $C^1$.  Then
\[
  \oint_{\partial D}(P\,dx+Q\,dy)
  = \iint_D\Bigl(\frac{\partial Q}{\partial x}
    -\frac{\partial P}{\partial y}\Bigr)\,dA.
\]
\end{theorem}

\begin{remark}[Proof idea]
For a Type~I domain $D=\{(x,y):a\le x\le b,\;\varphi_1(x)\le y\le\varphi_2(x)\}$,
one shows separately that
\[
  \oint_{\partial D}P\,dx = -\iint_D\frac{\partial P}{\partial y}\,dA
  \qquad\text{and}\qquad
  \oint_{\partial D}Q\,dy = \iint_D\frac{\partial Q}{\partial x}\,dA,
\]
by direct computation using Fubini's theorem.  For the first identity, the
boundary integral over $\partial D$ splits into a ``bottom'' piece
$y=\varphi_1(x)$ (traversed left-to-right) and a ``top'' piece
$y=\varphi_2(x)$ (right-to-left), and one verifies:
\[
  \oint_{\partial D}P\,dx
  = \int_a^b P(x,\varphi_1(x))\,dx - \int_a^b P(x,\varphi_2(x))\,dx
  = -\int_a^b\int_{\varphi_1(x)}^{\varphi_2(x)}
    \frac{\partial P}{\partial y}\,dy\,dx.
\]
The argument for $Q$ is analogous using a Type~II decomposition.  General
domains are handled by decomposing into regions of these types.
\end{remark}

\begin{center}
\begin{tikzpicture}[scale=1.3, decoration={markings,
  mark=between positions 0.05 and 0.95 step 0.1
  with {\arrow{Stealth[length=2.5mm]}}}]
  % A closed region D
  \draw[thick, blue!70!black, postaction={decorate}]
    (0,0) .. controls (1.5,2.5) and (3,2.8) .. (4,1.5)
    .. controls (3.5,0) and (1.5,-0.5) .. (0,0);
  \fill[blue!8]
    (0,0) .. controls (1.5,2.5) and (3,2.8) .. (4,1.5)
    .. controls (3.5,0) and (1.5,-0.5) .. (0,0);
  \node at (2,1.2) {$D$};
  \node[blue!70!black] at (4.5,2.0) {$\partial D$};
  % Double arrow equation
  \node at (2,-1.0) {$\displaystyle\oint_{\partial D}(P\,dx+Q\,dy)
    =\iint_D\Bigl(\frac{\partial Q}{\partial x}
    -\frac{\partial P}{\partial y}\Bigr)dA$};
\end{tikzpicture}
\captionof{figure}{Green's theorem: the circulation around $\partial D$
equals the double integral of the ``curl'' over $D$.}
\label{fig:green-theorem}
\end{center}

\begin{corollary}[Area via line integrals]
\label{cor:area-green}
\index{area!via Green's theorem}
Taking $P=-y/2$, $Q=x/2$, Green's theorem gives
\[
  \text{Area}(D) = \frac{1}{2}\oint_{\partial D}(x\,dy-y\,dx).
\]
\end{corollary}

\begin{example}[Area of an ellipse]
\label{ex:area-ellipse}
Parameterise $\partial D$ by $\gamma(t)=(a\cos t,b\sin t)$, $t\in[0,2\pi]$.
Then
\begin{align*}
  \text{Area} &= \frac{1}{2}\int_0^{2\pi}
    \bigl(a\cos t\cdot b\cos t - b\sin t\cdot(-a\sin t)\bigr)dt \\
  &= \frac{ab}{2}\int_0^{2\pi}(\cos^2 t+\sin^2 t)\,dt
  = \frac{ab}{2}\cdot 2\pi = \pi ab.
\end{align*}
\end{example}

%% ──────────────────────────────────────────────────────────
\section{A Closed Form That Is Not Exact}
\label{sec:closed-not-exact}
\index{closed form}\index{exact form}

\begin{example}[The angular form on $\mathbb{R}^2\setminus\{0\}$]
\label{ex:angular-form}
Consider the vector field on $U=\mathbb{R}^2\setminus\{(0,0)\}$:
\[
  F(x,y) = \Bigl(\frac{-y}{x^2+y^2},\;\frac{x}{x^2+y^2}\Bigr).
\]
Setting $P=\frac{-y}{x^2+y^2}$ and $Q=\frac{x}{x^2+y^2}$, one computes
\[
  \frac{\partial P}{\partial y}
  = \frac{-(x^2+y^2)+y\cdot 2y}{(x^2+y^2)^2}
  = \frac{y^2-x^2}{(x^2+y^2)^2}
  = \frac{\partial Q}{\partial x}.
\]
So the necessary condition $\partial P/\partial y=\partial Q/\partial x$ is
satisfied.  Yet $F$ is \emph{not} a gradient field on $U$, because
integrating around the unit circle $\gamma(t)=(\cos t,\sin t)$,
$t\in[0,2\pi]$:
\[
  \oint_\gamma F\cdot d\gamma
  = \int_0^{2\pi}\bigl((-\sin t)(-\sin t)+(\cos t)(\cos t)\bigr)dt
  = \int_0^{2\pi}1\,dt = 2\pi \ne 0.
\]
A gradient field would give zero around any closed curve, by
Theorem~\ref{thm:ftli}.  The obstruction is topological:
$U=\mathbb{R}^2\setminus\{0\}$ is not simply connected.
\end{example}

%% ──────────────────────────────────────────────────────────
\section{Exercises}
\label{sec:ex-ch11}

\begin{exercise}
Compute the arc length of the parabola $\gamma(t)=(t,t^2)$ for
$t\in[0,1]$.
\end{exercise}

\begin{exercise}
Compute the arc length of one arch of the cycloid
$\gamma(t)=(t-\sin t,1-\cos t)$ for $t\in[0,2\pi]$.
\end{exercise}

\begin{exercise}
Find the arc length parameterisation of $\gamma(t)=(3\cos t,3\sin t,4t)$,
$t\ge 0$.
\end{exercise}

\begin{exercise}
Compute $\int_\gamma (x^2+y^2)\,ds$ where $\gamma$ is the circle of radius
$R$ centred at the origin.
\end{exercise}

\begin{exercise}
Compute $\int_\gamma F\cdot d\gamma$ where $F(x,y)=(y,x)$ and
$\gamma(t)=(t,t^2)$, $t\in[0,1]$.
\end{exercise}

\begin{exercise}
Show that $F(x,y)=(2xy+3,x^2-4y)$ is conservative.  Find a potential $V$
and use it to compute $\int_\gamma F\cdot d\gamma$ from $(0,0)$ to $(1,2)$.
\end{exercise}

\begin{exercise}
Use Green's theorem to evaluate
$\oint_\gamma(y^2\,dx+3xy\,dy)$
where $\gamma$ is the boundary of the square $[0,1]\times[0,1]$ traversed
counterclockwise.
\end{exercise}

\begin{exercise}
Use Green's theorem to compute the area enclosed by the astroid
$\gamma(t)=(\cos^3 t,\sin^3 t)$, $t\in[0,2\pi]$.
\end{exercise}

\begin{exercise}
Let $F(x,y,z)=(yz,xz,xy)$.  Show that $F=\nabla V$ for $V(x,y,z)=xyz$ and
compute $\int_\gamma F\cdot d\gamma$ along any path from $(1,0,0)$ to
$(0,1,1)$.
\end{exercise}

\begin{exercise}
Let $F(x,y)=\bigl(\frac{-y}{x^2+y^2},\frac{x}{x^2+y^2}\bigr)$ on
$U=\mathbb{R}^2\setminus\{0\}$.  Compute $\int_\gamma F\cdot d\gamma$ where
$\gamma$ traverses the circle of radius~$2$ centred at the origin once
counterclockwise.  Does the result change if $\gamma$ winds twice around the
origin?
\end{exercise}

\begin{exercise}[$\star\star$]
Let $D\subseteq\mathbb{R}^2$ be a simply connected open set and
$F=(P,Q)\in C^1(D;\mathbb{R}^2)$ with
$\frac{\partial P}{\partial y}=\frac{\partial Q}{\partial x}$ everywhere in
$D$.  Prove that $F$ is conservative, i.e.\ construct a potential
$V(x,y)=\int_{(x_0,y_0)}^{(x,y)}F\cdot d\gamma$ and show it is well
defined.
\end{exercise}

\subsection*{Chapter Summary}

\begin{itemize}
  \item A parametric curve $\gamma\colon[a,b]\to\mathbb{R}^n$ is regular if
    $\gamma'(t)\ne 0$; its arc length is $L=\int_a^b\|\gamma'(t)\|\,dt$.
  \item The scalar line integral $\int_\gamma f\,ds$ integrates a function
    with respect to arc length; the vector line integral
    $\int_\gamma F\cdot d\gamma$ computes work.
  \item For a gradient field $F=\nabla V$, the line integral depends only on
    the endpoints: $\int_\gamma F\cdot d\gamma=V(\text{end})-V(\text{start})$.
  \item Green's theorem connects line integrals around a closed curve to
    double integrals over the enclosed region.
  \item A closed 1-form need not be exact on a non-simply-connected domain.
\end{itemize}

%% ──────────────────────────────────────────────────────────
\section*{Concluding Remarks: Where to Go Next}

The material in this course provides the foundation for several beautiful and
deep areas of mathematics.

\begin{itemize}
  \item \textbf{Measure theory and Lebesgue integration.}\; The Riemann
    integral, while powerful, has significant limitations (for instance, the
    limit of a pointwise convergent sequence of Riemann integrable functions
    need not be Riemann integrable).  Lebesgue's theory provides a far more
    flexible framework, with powerful convergence theorems (monotone
    convergence, dominated convergence) that are indispensable in
    probability, functional analysis, and PDE theory.
  \item \textbf{Complex analysis.}\; Functions of a complex variable
    $f\colon\mathbb{C}\to\mathbb{C}$ that are differentiable (holomorphic)
    possess extraordinary properties: they are automatically $C^\infty$,
    analytic, and satisfy Cauchy's integral formula.  The residue theorem is
    a far-reaching generalisation of the techniques used to compute
    $\int e^{-x^2}dx$.
  \item \textbf{Differential geometry and manifolds.}\; Green's theorem is a
    special case of the general Stokes theorem on manifolds, which unifies
    Green's, Gauss's, and Stokes's classical theorems.  The language of
    differential forms provides the natural setting for integration on curved
    spaces.
  \item \textbf{Partial differential equations.}\; Many of the tools
    developed here --- multivariable derivatives, Green's theorem,
    change of variables --- are the starting point for the study of the
    heat equation, wave equation, and Laplace's equation.
\end{itemize}


%% ══════════════════════════════════════════════════════════
%%  BIBLIOGRAPHY
%% ══════════════════════════════════════════════════════════

\begin{thebibliography}{99}

\bibitem{rudin}
W.~Rudin,
\emph{Principles of Mathematical Analysis},
3rd ed., McGraw-Hill, 1976.
\index{Rudin}

\bibitem{zuily-queffelec}
C.~Zuily and H.~Queff\'elec,
\emph{Analyse pour l'agr\'egation},
4th ed., Dunod, 2013.
\index{Zuily}\index{Queff\'elec}

\bibitem{ramis-deschamps-odoux}
E.~Ramis, C.~Deschamps, and J.~Odoux,
\emph{Cours de math\'ematiques sp\'eciales~-- Analyse},
Masson, 1991.
\index{Ramis}\index{Deschamps}\index{Odoux}

\bibitem{tao2}
T.~Tao,
\emph{Analysis~II},
3rd ed., Texts and Readings in Mathematics, Hindustan Book Agency, 2014.
\index{Tao}

\bibitem{pompeiani-tarrago}
I.~Pompe\"iani and T.~Tarrago,
\emph{Analyse: cours et exercices},
Ellipses, 2008.
\index{Pompe\"iani}\index{Tarrago}

\bibitem{spivak-manifolds}
M.~Spivak,
\emph{Calculus on Manifolds},
W.~A.~Benjamin, 1965.
\index{Spivak}

\end{thebibliography}


%% ══════════════════════════════════════════════════════════
%%  INDEX
%% ══════════════════════════════════════════════════════════

\printindex

\end{document}
